% Supplementary material of the AnytimeSwarm-CC manuscript (material moved out of the main text when it was shortened).
% Build after main: pdflatex supplementary && bibtex supplementary && pdflatex supplementary && pdflatex supplementary
% Every number comes from the generated tables/numbers*.tex files; none is typed by hand.
\documentclass[preprint,11pt]{elsarticle}
\usepackage[T1]{fontenc}
\usepackage{newtxtext}
\usepackage[varvw]{newtxmath}
\usepackage{graphicx}
\usepackage{booktabs}
\usepackage{amsmath}
\usepackage{xcolor}
\usepackage{xr-hyper}
\usepackage[hidelinks]{hyperref}
\usepackage[margin=2.5cm]{geometry}
\externaldocument{main}
\biboptions{sort&compress}
\newtheorem{prop}{Proposition}
\renewcommand{\theprop}{S\arabic{prop}}
\newproof{proof}{Proof}
\newcommand{\todo}[1]{\textcolor{red}{[TODO: #1]}}
\IfFileExists{tables/static_numbers.tex}{% copied from the static benchmark (conference pipeline); generated by experiments/08_make_tables.py --split test (D:\paper\TAIMMA_2027\AnytimeSwarm\configs\frozen_2026-10-24_test2.yaml) -- do not edit
\expandafter\def\csname vmu\endcsname{10^{-5}}
\expandafter\def\csname vMatchedB\endcsname{40}
\expandafter\def\csname vPSFixedSmall\endcsname{10}
\expandafter\def\csname vPSFixedMedium\endcsname{30}
\expandafter\def\csname vPSFixedLarge\endcsname{75}
\expandafter\def\csname vPSStagSmall\endcsname{30}
\expandafter\def\csname vPSStagMedium\endcsname{20}
\expandafter\def\csname vPSStagLarge\endcsname{50}
\expandafter\def\csname vNsnapTest\endcsname{1562}
\expandafter\def\csname vB\endcsname{250}
\expandafter\def\csname vASRegret\endcsname{0.08}
\expandafter\def\csname vASMs\endcsname{39.3}
\expandafter\def\csname vASUsed\endcsname{16}
\expandafter\def\csname vFixedRegret\endcsname{0.06}
\expandafter\def\csname vFixedMs\endcsname{250.1}
\expandafter\def\csname vFixedUsed\endcsname{100}
\expandafter\def\csname vStagRegret\endcsname{0.11}
\expandafter\def\csname vStagMs\endcsname{56.8}
\expandafter\def\csname vStagUsed\endcsname{23}
\expandafter\def\csname vRepairRegret\endcsname{0.98}
\expandafter\def\csname vRepairMs\endcsname{0.5}
\expandafter\def\csname vRepairUsed\endcsname{0}
\expandafter\def\csname vOracleRegret\endcsname{0.07}
\expandafter\def\csname vOracleMs\endcsname{12.5}
\expandafter\def\csname vOracleUsed\endcsname{5}
\expandafter\def\csname vdRegretFixed\endcsname{0.021}
\expandafter\def\csname vdRegretFixedCI\endcsname{[0.014,\,0.029]}
\expandafter\def\csname vdMsFixed\endcsname{-210.7}
\expandafter\def\csname vdMsFixedCI\endcsname{[-215.2,\,-206.5]}
\expandafter\def\csname vdRegretStag\endcsname{-0.025}
\expandafter\def\csname vdRegretStagCI\endcsname{[-0.038,\,-0.013]}
\expandafter\def\csname vdMsStag\endcsname{-17.5}
\expandafter\def\csname vdMsStagCI\endcsname{[-21.0,\,-14.0]}
\expandafter\def\csname vSavedPct\endcsname{84}
\expandafter\def\csname vKmax\endcsname{50}
\expandafter\def\csname vASEffK\endcsname{0.09}
\expandafter\def\csname vFixedEffK\endcsname{0.91}
\expandafter\def\csname vStagEffK\endcsname{0.21}
\expandafter\def\csname vNoDelayEffK\endcsname{0.13}
\expandafter\def\csname vOracleEffK\endcsname{0.07}
\expandafter\def\csname vdEffFixed\endcsname{-0.804}
\expandafter\def\csname vdEffFixedCI\endcsname{[-1.031,\,-0.597]}
\expandafter\def\csname vdEffStag\endcsname{-0.125}
\expandafter\def\csname vdEffStagCI\endcsname{[-0.219,\,-0.059]}
\expandafter\def\csname vdEffNoDelay\endcsname{-0.041}
\expandafter\def\csname vdEffNoDelayCI\endcsname{[-0.068,\,-0.016]}
\expandafter\def\csname vOnlineRegret\endcsname{0.31}
\expandafter\def\csname vOnlineMs\endcsname{7.2}
\expandafter\def\csname vdRegretOnline\endcsname{-0.226}
\expandafter\def\csname vdRegretOnlineCI\endcsname{[-0.281,\,-0.174]}
\expandafter\def\csname vBsmall\endcsname{20}
\expandafter\def\csname vASRegretSmallB\endcsname{0.20}
\expandafter\def\csname vNoFFRRegretSmallB\endcsname{0.39}
\expandafter\def\csname vdRegretNoFFRSmallB\endcsname{-0.193}
\expandafter\def\csname vdRegretNoFFRSmallBCI\endcsname{[-0.277,\,-0.122]}
\expandafter\def\csname vdRegretNoFFR\endcsname{-0.017}
\expandafter\def\csname vdRegretNoFFRCI\endcsname{[-0.034,\,-0.001]}
\expandafter\def\csname vdRegretNoFFRAbs\endcsname{0.017}
\expandafter\def\csname vTffLargePninetyfive\endcsname{3.3}
\expandafter\def\csname vTffLargeMedian\endcsname{0.9}
\expandafter\def\csname vOrtN\endcsname{1025}
\expandafter\def\csname vOrtWithin\endcsname{59}
\expandafter\def\csname vOrtBetter\endcsname{0}
\expandafter\def\csname vFrontEffFixedDiff\endcsname{-0.139}
\expandafter\def\csname vFrontEffFixedCI\endcsname{[-0.229,\,-0.071]}
\expandafter\def\csname vFrontEffStagDiff\endcsname{-0.131}
\expandafter\def\csname vFrontEffStagCI\endcsname{[-0.224,\,-0.067]}
\expandafter\def\csname vFrontEffNoDelayDiff\endcsname{-0.046}
\expandafter\def\csname vFrontEffNoDelayCI\endcsname{[-0.074,\,-0.021]}
\expandafter\def\csname vFrontEffFixedPSDiff\endcsname{-0.083}
\expandafter\def\csname vFrontEffFixedPSCI\endcsname{[-0.112,\,-0.055]}
\expandafter\def\csname vFrontEffStagPSDiff\endcsname{-0.120}
\expandafter\def\csname vFrontEffStagPSCI\endcsname{[-0.211,\,-0.058]}
\expandafter\def\csname vFrontFixedParam\endcsname{50}
\expandafter\def\csname vFrontFixedMs\endcsname{50.0}
\expandafter\def\csname vFrontOursMs\endcsname{41.5}
\expandafter\def\csname vFrontFixedDiff\endcsname{-0.045}
\expandafter\def\csname vFrontFixedCI\endcsname{[-0.062,\,-0.028]}
\expandafter\def\csname vFrontFixedBase\endcsname{0.127}
\expandafter\def\csname vFrontOursRegret\endcsname{0.083}
\expandafter\def\csname vFrontFixedRed\endcsname{35}
\expandafter\def\csname vFrontFixedMoreMs\endcsname{20}
\expandafter\def\csname vFrontStagParam\endcsname{50}
\expandafter\def\csname vFrontStagMs\endcsname{56.5}
\expandafter\def\csname vFrontStagDiff\endcsname{-0.029}
\expandafter\def\csname vFrontStagCI\endcsname{[-0.043,\,-0.016]}
\expandafter\def\csname vFrontStagBase\endcsname{0.112}
\expandafter\def\csname vFrontStagRed\endcsname{26}
\expandafter\def\csname vFrontStagMoreMs\endcsname{36}
\expandafter\def\csname vFrontFixedPSParam\endcsname{0}
\expandafter\def\csname vFrontFixedPSMs\endcsname{38.3}
\expandafter\def\csname vFrontFixedPSDiff\endcsname{-0.047}
\expandafter\def\csname vFrontFixedPSCI\endcsname{[-0.064,\,-0.031]}
\expandafter\def\csname vFrontFixedPSBase\endcsname{0.130}
\expandafter\def\csname vFrontFixedPSRed\endcsname{36}
\expandafter\def\csname vFrontFixedPSMoreMs\endcsname{-8}
\expandafter\def\csname vFrontStagPSParam\endcsname{0}
\expandafter\def\csname vFrontStagPSMs\endcsname{39.2}
\expandafter\def\csname vFrontStagPSDiff\endcsname{-0.039}
\expandafter\def\csname vFrontStagPSCI\endcsname{[-0.055,\,-0.026]}
\expandafter\def\csname vFrontStagPSBase\endcsname{0.122}
\expandafter\def\csname vFrontStagPSRed\endcsname{32}
\expandafter\def\csname vFrontStagPSMoreMs\endcsname{-5}
\expandafter\def\csname vCLASUtil1\endcsname{88.9}
\expandafter\def\csname vCLFixedUtil1\endcsname{89.0}
\expandafter\def\csname vCLStagUtil1\endcsname{88.9}
\expandafter\def\csname vCLFixedMatchedUtil1\endcsname{88.6}
\expandafter\def\csname vCLRepairUtil1\endcsname{84.0}
\expandafter\def\csname vCLNoDelayUtil1\endcsname{88.9}
\expandafter\def\csname vCLFixedPSUtil1\endcsname{88.8}
\expandafter\def\csname vCLStagPSUtil1\endcsname{88.9}
\expandafter\def\csname vCLASDrop1\endcsname{0.03}
\expandafter\def\csname vCLASPlanMs1\endcsname{46.2}
\expandafter\def\csname vCLFixedDrop1\endcsname{0.10}
\expandafter\def\csname vCLFixedPlanMs1\endcsname{250.1}
\expandafter\def\csname vCLStagDrop1\endcsname{0.05}
\expandafter\def\csname vCLStagPlanMs1\endcsname{62.7}
\expandafter\def\csname vCLFixedMatchedDrop1\endcsname{0.00}
\expandafter\def\csname vCLFixedMatchedPlanMs1\endcsname{40.1}
\expandafter\def\csname vCLRepairDrop1\endcsname{0.00}
\expandafter\def\csname vCLRepairPlanMs1\endcsname{1.1}
\expandafter\def\csname vCLNoDelayDrop1\endcsname{0.06}
\expandafter\def\csname vCLNoDelayPlanMs1\endcsname{47.1}
\expandafter\def\csname vCLFixedPSDrop1\endcsname{0.00}
\expandafter\def\csname vCLFixedPSPlanMs1\endcsname{38.4}
\expandafter\def\csname vCLStagPSDrop1\endcsname{0.04}
\expandafter\def\csname vCLStagPSPlanMs1\endcsname{43.7}
\expandafter\def\csname vCLdUtilFixed1\endcsname{-0.07}
\expandafter\def\csname vCLdUtilFixed1CI\endcsname{[-0.35,\,0.20]}
\expandafter\def\csname vCLdUtilStag1\endcsname{0.08}
\expandafter\def\csname vCLdUtilStag1CI\endcsname{[-0.17,\,0.29]}
\expandafter\def\csname vCLdUtilFixedMatched1\endcsname{0.34}
\expandafter\def\csname vCLdUtilFixedMatched1CI\endcsname{[0.10,\,0.55]}
\expandafter\def\csname vCLdUtilNoDelay1\endcsname{0.00}
\expandafter\def\csname vCLdUtilNoDelay1CI\endcsname{[-0.04,\,0.05]}
\expandafter\def\csname vCLdUtilFixedPS1\endcsname{0.09}
\expandafter\def\csname vCLdUtilFixedPS1CI\endcsname{[-0.20,\,0.34]}
\expandafter\def\csname vCLdUtilStagPS1\endcsname{0.01}
\expandafter\def\csname vCLdUtilStagPS1CI\endcsname{[-0.23,\,0.21]}
\expandafter\def\csname vCLdUtilFixedOneSmall\endcsname{-0.41}
\expandafter\def\csname vCLdUtilFixedOneSmallCI\endcsname{[-1.13,\,0.27]}
\expandafter\def\csname vCLdUtilStagPSOneSmall\endcsname{-0.54}
\expandafter\def\csname vCLdUtilStagPSOneSmallCI\endcsname{[-1.05,\,-0.10]}
\expandafter\def\csname vCLdUtilFixedPSOneSmall\endcsname{-0.43}
\expandafter\def\csname vCLdUtilFixedPSOneSmallCI\endcsname{[-1.15,\,0.17]}
\expandafter\def\csname vCLdUtilFixedOneMedium\endcsname{0.28}
\expandafter\def\csname vCLdUtilFixedOneMediumCI\endcsname{[-0.20,\,0.76]}
\expandafter\def\csname vCLdUtilStagPSOneMedium\endcsname{0.28}
\expandafter\def\csname vCLdUtilStagPSOneMediumCI\endcsname{[-0.02,\,0.60]}
\expandafter\def\csname vCLdUtilFixedPSOneMedium\endcsname{0.27}
\expandafter\def\csname vCLdUtilFixedPSOneMediumCI\endcsname{[-0.10,\,0.65]}
\expandafter\def\csname vCLdUtilFixedOneLarge\endcsname{-0.07}
\expandafter\def\csname vCLdUtilFixedOneLargeCI\endcsname{[-0.21,\,0.07]}
\expandafter\def\csname vCLdUtilStagPSOneLarge\endcsname{0.27}
\expandafter\def\csname vCLdUtilStagPSOneLargeCI\endcsname{[0.06,\,0.49]}
\expandafter\def\csname vCLdUtilFixedPSOneLarge\endcsname{0.44}
\expandafter\def\csname vCLdUtilFixedPSOneLargeCI\endcsname{[0.24,\,0.65]}
\expandafter\def\csname vCLASUtil2\endcsname{88.9}
\expandafter\def\csname vCLFixedUtil2\endcsname{89.0}
\expandafter\def\csname vCLStagUtil2\endcsname{88.9}
\expandafter\def\csname vCLFixedMatchedUtil2\endcsname{88.6}
\expandafter\def\csname vCLRepairUtil2\endcsname{84.0}
\expandafter\def\csname vCLNoDelayUtil2\endcsname{88.9}
\expandafter\def\csname vCLFixedPSUtil2\endcsname{88.8}
\expandafter\def\csname vCLStagPSUtil2\endcsname{88.9}
\expandafter\def\csname vCLASDrop2\endcsname{0.03}
\expandafter\def\csname vCLASPlanMs2\endcsname{45.1}
\expandafter\def\csname vCLFixedDrop2\endcsname{0.18}
\expandafter\def\csname vCLFixedPlanMs2\endcsname{250.1}
\expandafter\def\csname vCLStagDrop2\endcsname{0.09}
\expandafter\def\csname vCLStagPlanMs2\endcsname{62.7}
\expandafter\def\csname vCLFixedMatchedDrop2\endcsname{0.01}
\expandafter\def\csname vCLFixedMatchedPlanMs2\endcsname{40.1}
\expandafter\def\csname vCLRepairDrop2\endcsname{0.01}
\expandafter\def\csname vCLRepairPlanMs2\endcsname{1.1}
\expandafter\def\csname vCLNoDelayDrop2\endcsname{0.12}
\expandafter\def\csname vCLNoDelayPlanMs2\endcsname{47.0}
\expandafter\def\csname vCLFixedPSDrop2\endcsname{0.02}
\expandafter\def\csname vCLFixedPSPlanMs2\endcsname{38.4}
\expandafter\def\csname vCLStagPSDrop2\endcsname{0.08}
\expandafter\def\csname vCLStagPSPlanMs2\endcsname{43.6}
\expandafter\def\csname vCLASUtil5\endcsname{88.9}
\expandafter\def\csname vCLFixedUtil5\endcsname{88.7}
\expandafter\def\csname vCLStagUtil5\endcsname{88.8}
\expandafter\def\csname vCLFixedMatchedUtil5\endcsname{88.5}
\expandafter\def\csname vCLRepairUtil5\endcsname{84.0}
\expandafter\def\csname vCLNoDelayUtil5\endcsname{88.8}
\expandafter\def\csname vCLFixedPSUtil5\endcsname{88.8}
\expandafter\def\csname vCLStagPSUtil5\endcsname{88.8}
\expandafter\def\csname vCLASDrop5\endcsname{0.10}
\expandafter\def\csname vCLASPlanMs5\endcsname{43.3}
\expandafter\def\csname vCLFixedDrop5\endcsname{0.47}
\expandafter\def\csname vCLFixedPlanMs5\endcsname{250.1}
\expandafter\def\csname vCLStagDrop5\endcsname{0.22}
\expandafter\def\csname vCLStagPlanMs5\endcsname{62.6}
\expandafter\def\csname vCLFixedMatchedDrop5\endcsname{0.06}
\expandafter\def\csname vCLFixedMatchedPlanMs5\endcsname{40.1}
\expandafter\def\csname vCLRepairDrop5\endcsname{0.01}
\expandafter\def\csname vCLRepairPlanMs5\endcsname{1.1}
\expandafter\def\csname vCLNoDelayDrop5\endcsname{0.36}
\expandafter\def\csname vCLNoDelayPlanMs5\endcsname{47.2}
\expandafter\def\csname vCLFixedPSDrop5\endcsname{0.11}
\expandafter\def\csname vCLFixedPSPlanMs5\endcsname{38.4}
\expandafter\def\csname vCLStagPSDrop5\endcsname{0.19}
\expandafter\def\csname vCLStagPSPlanMs5\endcsname{43.6}
\expandafter\def\csname vCLASUtil10\endcsname{88.8}
\expandafter\def\csname vCLFixedUtil10\endcsname{88.2}
\expandafter\def\csname vCLStagUtil10\endcsname{88.6}
\expandafter\def\csname vCLFixedMatchedUtil10\endcsname{88.3}
\expandafter\def\csname vCLRepairUtil10\endcsname{84.0}
\expandafter\def\csname vCLNoDelayUtil10\endcsname{88.6}
\expandafter\def\csname vCLFixedPSUtil10\endcsname{88.6}
\expandafter\def\csname vCLStagPSUtil10\endcsname{88.7}
\expandafter\def\csname vCLASDrop10\endcsname{0.17}
\expandafter\def\csname vCLASPlanMs10\endcsname{41.3}
\expandafter\def\csname vCLFixedDrop10\endcsname{1.01}
\expandafter\def\csname vCLFixedPlanMs10\endcsname{250.1}
\expandafter\def\csname vCLStagDrop10\endcsname{0.44}
\expandafter\def\csname vCLStagPlanMs10\endcsname{63.3}
\expandafter\def\csname vCLFixedMatchedDrop10\endcsname{0.16}
\expandafter\def\csname vCLFixedMatchedPlanMs10\endcsname{40.1}
\expandafter\def\csname vCLRepairDrop10\endcsname{0.02}
\expandafter\def\csname vCLRepairPlanMs10\endcsname{1.1}
\expandafter\def\csname vCLNoDelayDrop10\endcsname{0.71}
\expandafter\def\csname vCLNoDelayPlanMs10\endcsname{47.7}
\expandafter\def\csname vCLFixedPSDrop10\endcsname{0.23}
\expandafter\def\csname vCLFixedPSPlanMs10\endcsname{38.4}
\expandafter\def\csname vCLStagPSDrop10\endcsname{0.39}
\expandafter\def\csname vCLStagPSPlanMs10\endcsname{44.2}
\expandafter\def\csname vCLASUtil20\endcsname{88.6}
\expandafter\def\csname vCLFixedUtil20\endcsname{87.7}
\expandafter\def\csname vCLStagUtil20\endcsname{88.4}
\expandafter\def\csname vCLFixedMatchedUtil20\endcsname{88.2}
\expandafter\def\csname vCLRepairUtil20\endcsname{84.0}
\expandafter\def\csname vCLNoDelayUtil20\endcsname{88.3}
\expandafter\def\csname vCLFixedPSUtil20\endcsname{88.5}
\expandafter\def\csname vCLStagPSUtil20\endcsname{88.4}
\expandafter\def\csname vCLASDrop20\endcsname{0.20}
\expandafter\def\csname vCLASPlanMs20\endcsname{39.0}
\expandafter\def\csname vCLFixedDrop20\endcsname{1.84}
\expandafter\def\csname vCLFixedPlanMs20\endcsname{250.1}
\expandafter\def\csname vCLStagDrop20\endcsname{0.88}
\expandafter\def\csname vCLStagPlanMs20\endcsname{63.2}
\expandafter\def\csname vCLFixedMatchedDrop20\endcsname{0.30}
\expandafter\def\csname vCLFixedMatchedPlanMs20\endcsname{40.1}
\expandafter\def\csname vCLRepairDrop20\endcsname{0.06}
\expandafter\def\csname vCLRepairPlanMs20\endcsname{1.1}
\expandafter\def\csname vCLNoDelayDrop20\endcsname{1.25}
\expandafter\def\csname vCLNoDelayPlanMs20\endcsname{47.4}
\expandafter\def\csname vCLFixedPSDrop20\endcsname{0.47}
\expandafter\def\csname vCLFixedPSPlanMs20\endcsname{38.4}
\expandafter\def\csname vCLStagPSDrop20\endcsname{0.75}
\expandafter\def\csname vCLStagPSPlanMs20\endcsname{43.8}
\expandafter\def\csname vCLASUtil50\endcsname{88.3}
\expandafter\def\csname vCLFixedUtil50\endcsname{85.5}
\expandafter\def\csname vCLStagUtil50\endcsname{87.3}
\expandafter\def\csname vCLFixedMatchedUtil50\endcsname{87.9}
\expandafter\def\csname vCLRepairUtil50\endcsname{84.0}
\expandafter\def\csname vCLNoDelayUtil50\endcsname{87.2}
\expandafter\def\csname vCLFixedPSUtil50\endcsname{88.0}
\expandafter\def\csname vCLStagPSUtil50\endcsname{87.6}
\expandafter\def\csname vCLASDrop50\endcsname{0.39}
\expandafter\def\csname vCLASPlanMs50\endcsname{35.8}
\expandafter\def\csname vCLFixedDrop50\endcsname{4.47}
\expandafter\def\csname vCLFixedPlanMs50\endcsname{250.1}
\expandafter\def\csname vCLStagDrop50\endcsname{2.13}
\expandafter\def\csname vCLStagPlanMs50\endcsname{63.5}
\expandafter\def\csname vCLFixedMatchedDrop50\endcsname{0.73}
\expandafter\def\csname vCLFixedMatchedPlanMs50\endcsname{40.1}
\expandafter\def\csname vCLRepairDrop50\endcsname{0.19}
\expandafter\def\csname vCLRepairPlanMs50\endcsname{1.1}
\expandafter\def\csname vCLNoDelayDrop50\endcsname{2.92}
\expandafter\def\csname vCLNoDelayPlanMs50\endcsname{48.7}
\expandafter\def\csname vCLFixedPSDrop50\endcsname{1.07}
\expandafter\def\csname vCLFixedPSPlanMs50\endcsname{38.4}
\expandafter\def\csname vCLStagPSDrop50\endcsname{1.82}
\expandafter\def\csname vCLStagPSPlanMs50\endcsname{44.0}
\expandafter\def\csname vCLdUtilFixed\endcsname{2.81}
\expandafter\def\csname vCLdUtilFixedCI\endcsname{[2.31,\,3.30]}
\expandafter\def\csname vCLdUtilStag\endcsname{1.01}
\expandafter\def\csname vCLdUtilStagCI\endcsname{[0.67,\,1.36]}
\expandafter\def\csname vCLdUtilFixedMatched\endcsname{0.40}
\expandafter\def\csname vCLdUtilFixedMatchedCI\endcsname{[0.12,\,0.68]}
\expandafter\def\csname vCLdUtilNoDelay\endcsname{1.09}
\expandafter\def\csname vCLdUtilNoDelayCI\endcsname{[0.83,\,1.36]}
\expandafter\def\csname vCLdUtilFixedPS\endcsname{0.26}
\expandafter\def\csname vCLdUtilFixedPSCI\endcsname{[-0.01,\,0.56]}
\expandafter\def\csname vCLdUtilStagPS\endcsname{0.65}
\expandafter\def\csname vCLdUtilStagPSCI\endcsname{[0.37,\,0.92]}
\expandafter\def\csname vNMissions\endcsname{270}
\expandafter\def\csname vNKappas\endcsname{6}
\expandafter\def\csname vKappaWin\endcsname{10}
\expandafter\def\csname vKCdiff1\endcsname{-0.07}
\expandafter\def\csname vKCdiff1CI\endcsname{[-0.35,\,0.22]}
\expandafter\def\csname vKCdiff2\endcsname{-0.10}
\expandafter\def\csname vKCdiff2CI\endcsname{[-0.37,\,0.18]}
\expandafter\def\csname vKCdiff5\endcsname{0.18}
\expandafter\def\csname vKCdiff5CI\endcsname{[-0.14,\,0.50]}
\expandafter\def\csname vKCdiff10\endcsname{0.56}
\expandafter\def\csname vKCdiff10CI\endcsname{[0.21,\,0.91]}
\expandafter\def\csname vKCdiff20\endcsname{0.98}
\expandafter\def\csname vKCdiff20CI\endcsname{[0.62,\,1.35]}
\expandafter\def\csname vKCdiff50\endcsname{2.81}
\expandafter\def\csname vKCdiff50CI\endcsname{[2.34,\,3.30]}
\expandafter\def\csname vEqRegretAS\endcsname{0.083}
\expandafter\def\csname vEqEffAS\endcsname{0.095}
\expandafter\def\csname vEqRegretBaseMin\endcsname{0.112}
\expandafter\def\csname vEqRegretBaseMax\endcsname{0.130}
\expandafter\def\csname vEqEffBaseMin\endcsname{0.177}
\expandafter\def\csname vEqEffBaseMax\endcsname{0.233}
\expandafter\def\csname vEqRegretRedMin\endcsname{26}
\expandafter\def\csname vEqRegretRedMax\endcsname{36}
\expandafter\def\csname vNeedFixedMs\endcsname{141}
\expandafter\def\csname vSpeedFixed\endcsname{3.4}
\expandafter\def\csname vNeedStagMs\endcsname{98}
\expandafter\def\csname vSpeedStag\endcsname{2.4}
\expandafter\def\csname vRedSmallMin\endcsname{-0}
\expandafter\def\csname vRedSmallMax\endcsname{82}
\expandafter\def\csname vRedMediumMin\endcsname{-7}
\expandafter\def\csname vRedMediumMax\endcsname{46}
\expandafter\def\csname vRedLargeMin\endcsname{29}
\expandafter\def\csname vRedLargeMax\endcsname{39}
\expandafter\def\csname vDropRedFixed\endcsname{91}
\expandafter\def\csname vDropRedStag\endcsname{82}
\expandafter\def\csname vDropRedNoDelay\endcsname{87}
\expandafter\def\csname vPlanRatio1\endcsname{5.4}
\expandafter\def\csname vCLUtilBaseMinK\endcsname{85.5}
\expandafter\def\csname vCLUtilBaseMaxK\endcsname{88.0}
\expandafter\def\csname vAllocASNone\endcsname{18.2}
\expandafter\def\csname vAllocASHigh\endcsname{140.0}
\expandafter\def\csname vAllocASShare\endcsname{29}
\expandafter\def\csname vAllocFixedPSNone\endcsname{31.7}
\expandafter\def\csname vAllocFixedPSHigh\endcsname{43.5}
\expandafter\def\csname vAllocFixedPSShare\endcsname{54}
\expandafter\def\csname vAllocStagPSNone\endcsname{30.9}
\expandafter\def\csname vAllocStagPSHigh\endcsname{56.5}
\expandafter\def\csname vAllocStagPSShare\endcsname{52}
\expandafter\def\csname vAllocNoneShare\endcsname{67}
\expandafter\def\csname vAUCLearned\endcsname{0.926}
\expandafter\def\csname vAUCOnline\endcsname{0.757}
\expandafter\def\csname vPosRate\endcsname{1.3}
\expandafter\def\csname vASMsSmall\endcsname{18}
\expandafter\def\csname vASMsMedium\endcsname{24}
\expandafter\def\csname vASMsLarge\endcsname{74}
}{}
\IfFileExists{tables/static_numbers_XL.tex}{% copied from the static benchmark (conference pipeline); generated by experiments/08_make_tables.py --split test (D:\paper\TAIMMA_2027\AnytimeSwarm\configs\frozen_2026-10-24_xl.yaml) -- do not edit
\expandafter\def\csname vXLmu\endcsname{10^{-5}}
\expandafter\def\csname vXLMatchedB\endcsname{180}
\expandafter\def\csname vXLPSFixedXlarge\endcsname{150}
\expandafter\def\csname vXLPSStagXlarge\endcsname{100}
\expandafter\def\csname vXLNsnapTest\endcsname{534}
\expandafter\def\csname vXLB\endcsname{250}
\expandafter\def\csname vXLASRegret\endcsname{0.24}
\expandafter\def\csname vXLASMs\endcsname{122.6}
\expandafter\def\csname vXLASUsed\endcsname{49}
\expandafter\def\csname vXLFixedRegret\endcsname{0.23}
\expandafter\def\csname vXLFixedMs\endcsname{250.5}
\expandafter\def\csname vXLFixedUsed\endcsname{100}
\expandafter\def\csname vXLStagRegret\endcsname{0.29}
\expandafter\def\csname vXLStagMs\endcsname{73.0}
\expandafter\def\csname vXLStagUsed\endcsname{29}
\expandafter\def\csname vXLRepairRegret\endcsname{0.95}
\expandafter\def\csname vXLRepairMs\endcsname{5.3}
\expandafter\def\csname vXLRepairUsed\endcsname{2}
\expandafter\def\csname vXLOracleRegret\endcsname{0.24}
\expandafter\def\csname vXLOracleMs\endcsname{32.9}
\expandafter\def\csname vXLOracleUsed\endcsname{13}
\expandafter\def\csname vXLdRegretFixed\endcsname{0.013}
\expandafter\def\csname vXLdRegretFixedCI\endcsname{[0.009,\,0.018]}
\expandafter\def\csname vXLdMsFixed\endcsname{-126.8}
\expandafter\def\csname vXLdMsFixedCI\endcsname{[-135.4,\,-117.7]}
\expandafter\def\csname vXLdRegretStag\endcsname{-0.055}
\expandafter\def\csname vXLdRegretStagCI\endcsname{[-0.075,\,-0.037]}
\expandafter\def\csname vXLdMsStag\endcsname{50.1}
\expandafter\def\csname vXLdMsStagCI\endcsname{[42.7,\,57.6]}
\expandafter\def\csname vXLSavedPct\endcsname{51}
\expandafter\def\csname vXLKmax\endcsname{50}
\expandafter\def\csname vXLASEffK\endcsname{0.31}
\expandafter\def\csname vXLFixedEffK\endcsname{1.02}
\expandafter\def\csname vXLStagEffK\endcsname{0.46}
\expandafter\def\csname vXLNoDelayEffK\endcsname{0.58}
\expandafter\def\csname vXLOracleEffK\endcsname{0.26}
\expandafter\def\csname vXLdEffFixed\endcsname{-0.706}
\expandafter\def\csname vXLdEffFixedCI\endcsname{[-0.826,\,-0.596]}
\expandafter\def\csname vXLdEffStag\endcsname{-0.158}
\expandafter\def\csname vXLdEffStagCI\endcsname{[-0.213,\,-0.106]}
\expandafter\def\csname vXLdEffNoDelay\endcsname{-0.276}
\expandafter\def\csname vXLdEffNoDelayCI\endcsname{[-0.343,\,-0.211]}
\expandafter\def\csname vXLOnlineRegret\endcsname{0.57}
\expandafter\def\csname vXLOnlineMs\endcsname{18.4}
\expandafter\def\csname vXLdRegretOnline\endcsname{-0.336}
\expandafter\def\csname vXLdRegretOnlineCI\endcsname{[-0.428,\,-0.251]}
\expandafter\def\csname vXLBsmall\endcsname{20}
\expandafter\def\csname vXLASRegretSmallB\endcsname{0.57}
\expandafter\def\csname vXLNoFFRRegretSmallB\endcsname{1.42}
\expandafter\def\csname vXLdRegretNoFFRSmallB\endcsname{-0.868}
\expandafter\def\csname vXLdRegretNoFFRSmallBCI\endcsname{[-1.117,\,-0.611]}
\expandafter\def\csname vXLdRegretNoFFR\endcsname{-0.317}
\expandafter\def\csname vXLdRegretNoFFRCI\endcsname{[-0.463,\,-0.185]}
\expandafter\def\csname vXLTffLargePninetyfive\endcsname{14.7}
\expandafter\def\csname vXLTffLargeMedian\endcsname{4.0}
\expandafter\def\csname vXLFrontEffFixedDiff\endcsname{-0.357}
\expandafter\def\csname vXLFrontEffFixedCI\endcsname{[-0.443,\,-0.279]}
\expandafter\def\csname vXLFrontEffStagDiff\endcsname{-0.280}
\expandafter\def\csname vXLFrontEffStagCI\endcsname{[-0.352,\,-0.209]}
\expandafter\def\csname vXLFrontEffNoDelayDiff\endcsname{-0.340}
\expandafter\def\csname vXLFrontEffNoDelayCI\endcsname{[-0.421,\,-0.260]}
\expandafter\def\csname vXLFrontEffFixedPSDiff\endcsname{-0.357}
\expandafter\def\csname vXLFrontEffFixedPSCI\endcsname{[-0.443,\,-0.279]}
\expandafter\def\csname vXLFrontEffStagPSDiff\endcsname{-0.280}
\expandafter\def\csname vXLFrontEffStagPSCI\endcsname{[-0.352,\,-0.209]}
\expandafter\def\csname vXLFrontFixedParam\endcsname{150}
\expandafter\def\csname vXLFrontFixedMs\endcsname{150.0}
\expandafter\def\csname vXLFrontOursMs\endcsname{150.7}
\expandafter\def\csname vXLFrontFixedDiff\endcsname{-0.051}
\expandafter\def\csname vXLFrontFixedCI\endcsname{[-0.073,\,-0.031]}
\expandafter\def\csname vXLFrontFixedBase\endcsname{0.256}
\expandafter\def\csname vXLFrontOursRegret\endcsname{0.205}
\expandafter\def\csname vXLFrontFixedRed\endcsname{20}
\expandafter\def\csname vXLFrontFixedMoreMs\endcsname{-0}
\expandafter\def\csname vXLFrontStagParam\endcsname{200}
\expandafter\def\csname vXLFrontStagMs\endcsname{257.9}
\expandafter\def\csname vXLFrontStagDiff\endcsname{-0.000}
\expandafter\def\csname vXLFrontStagCI\endcsname{[-0.012,\,0.010]}
\expandafter\def\csname vXLFrontStagBase\endcsname{0.205}
\expandafter\def\csname vXLFrontStagRed\endcsname{0}
\expandafter\def\csname vXLFrontStagMoreMs\endcsname{71}
\expandafter\def\csname vXLFrontFixedPSParam\endcsname{0}
\expandafter\def\csname vXLFrontFixedPSMs\endcsname{150.0}
\expandafter\def\csname vXLFrontFixedPSDiff\endcsname{-0.051}
\expandafter\def\csname vXLFrontFixedPSCI\endcsname{[-0.073,\,-0.031]}
\expandafter\def\csname vXLFrontFixedPSBase\endcsname{0.256}
\expandafter\def\csname vXLFrontFixedPSRed\endcsname{20}
\expandafter\def\csname vXLFrontFixedPSMoreMs\endcsname{-0}
\expandafter\def\csname vXLFrontStagPSParam\endcsname{0}
\expandafter\def\csname vXLFrontStagPSMs\endcsname{134.8}
\expandafter\def\csname vXLFrontStagPSDiff\endcsname{-0.035}
\expandafter\def\csname vXLFrontStagPSCI\endcsname{[-0.053,\,-0.019]}
\expandafter\def\csname vXLFrontStagPSBase\endcsname{0.240}
\expandafter\def\csname vXLFrontStagPSRed\endcsname{15}
\expandafter\def\csname vXLFrontStagPSMoreMs\endcsname{-11}
\expandafter\def\csname vXLCLASUtil1\endcsname{93.1}
\expandafter\def\csname vXLCLFixedUtil1\endcsname{93.2}
\expandafter\def\csname vXLCLStagUtil1\endcsname{92.7}
\expandafter\def\csname vXLCLFixedMatchedUtil1\endcsname{92.8}
\expandafter\def\csname vXLCLRepairUtil1\endcsname{88.1}
\expandafter\def\csname vXLCLNoDelayUtil1\endcsname{93.2}
\expandafter\def\csname vXLCLFixedPSUtil1\endcsname{92.2}
\expandafter\def\csname vXLCLStagPSUtil1\endcsname{93.1}
\expandafter\def\csname vXLCLASDrop1\endcsname{0.32}
\expandafter\def\csname vXLCLASPlanMs1\endcsname{136.4}
\expandafter\def\csname vXLCLFixedDrop1\endcsname{0.56}
\expandafter\def\csname vXLCLFixedPlanMs1\endcsname{250.7}
\expandafter\def\csname vXLCLStagDrop1\endcsname{0.49}
\expandafter\def\csname vXLCLStagPlanMs1\endcsname{96.3}
\expandafter\def\csname vXLCLFixedMatchedDrop1\endcsname{0.38}
\expandafter\def\csname vXLCLFixedMatchedPlanMs1\endcsname{180.8}
\expandafter\def\csname vXLCLRepairDrop1\endcsname{0.21}
\expandafter\def\csname vXLCLRepairPlanMs1\endcsname{18.8}
\expandafter\def\csname vXLCLNoDelayDrop1\endcsname{0.56}
\expandafter\def\csname vXLCLNoDelayPlanMs1\endcsname{140.2}
\expandafter\def\csname vXLCLFixedPSDrop1\endcsname{0.43}
\expandafter\def\csname vXLCLFixedPSPlanMs1\endcsname{155.4}
\expandafter\def\csname vXLCLStagPSDrop1\endcsname{0.49}
\expandafter\def\csname vXLCLStagPSPlanMs1\endcsname{151.3}
\expandafter\def\csname vXLCLdUtilFixed1\endcsname{-0.10}
\expandafter\def\csname vXLCLdUtilFixed1CI\endcsname{[-0.21,\,0.01]}
\expandafter\def\csname vXLCLdUtilStag1\endcsname{0.38}
\expandafter\def\csname vXLCLdUtilStag1CI\endcsname{[0.20,\,0.60]}
\expandafter\def\csname vXLCLdUtilFixedMatched1\endcsname{0.25}
\expandafter\def\csname vXLCLdUtilFixedMatched1CI\endcsname{[0.10,\,0.42]}
\expandafter\def\csname vXLCLdUtilNoDelay1\endcsname{-0.08}
\expandafter\def\csname vXLCLdUtilNoDelay1CI\endcsname{[-0.19,\,0.02]}
\expandafter\def\csname vXLCLdUtilFixedPS1\endcsname{0.93}
\expandafter\def\csname vXLCLdUtilFixedPS1CI\endcsname{[0.65,\,1.21]}
\expandafter\def\csname vXLCLdUtilStagPS1\endcsname{0.03}
\expandafter\def\csname vXLCLdUtilStagPS1CI\endcsname{[-0.09,\,0.15]}
\expandafter\def\csname vXLCLdUtilFixedOneXlarge\endcsname{-0.10}
\expandafter\def\csname vXLCLdUtilFixedOneXlargeCI\endcsname{[-0.21,\,0.01]}
\expandafter\def\csname vXLCLdUtilStagPSOneXlarge\endcsname{0.03}
\expandafter\def\csname vXLCLdUtilStagPSOneXlargeCI\endcsname{[-0.09,\,0.15]}
\expandafter\def\csname vXLCLdUtilFixedPSOneXlarge\endcsname{0.93}
\expandafter\def\csname vXLCLdUtilFixedPSOneXlargeCI\endcsname{[0.65,\,1.21]}
\expandafter\def\csname vXLCLASUtil10\endcsname{92.5}
\expandafter\def\csname vXLCLFixedUtil10\endcsname{92.4}
\expandafter\def\csname vXLCLStagUtil10\endcsname{92.2}
\expandafter\def\csname vXLCLFixedMatchedUtil10\endcsname{92.3}
\expandafter\def\csname vXLCLRepairUtil10\endcsname{88.0}
\expandafter\def\csname vXLCLNoDelayUtil10\endcsname{92.5}
\expandafter\def\csname vXLCLFixedPSUtil10\endcsname{92.4}
\expandafter\def\csname vXLCLStagPSUtil10\endcsname{92.4}
\expandafter\def\csname vXLCLASDrop10\endcsname{2.24}
\expandafter\def\csname vXLCLASPlanMs10\endcsname{125.4}
\expandafter\def\csname vXLCLFixedDrop10\endcsname{4.58}
\expandafter\def\csname vXLCLFixedPlanMs10\endcsname{250.6}
\expandafter\def\csname vXLCLStagDrop10\endcsname{3.76}
\expandafter\def\csname vXLCLStagPlanMs10\endcsname{92.6}
\expandafter\def\csname vXLCLFixedMatchedDrop10\endcsname{3.53}
\expandafter\def\csname vXLCLFixedMatchedPlanMs10\endcsname{180.7}
\expandafter\def\csname vXLCLRepairDrop10\endcsname{2.52}
\expandafter\def\csname vXLCLRepairPlanMs10\endcsname{22.5}
\expandafter\def\csname vXLCLNoDelayDrop10\endcsname{4.36}
\expandafter\def\csname vXLCLNoDelayPlanMs10\endcsname{139.3}
\expandafter\def\csname vXLCLFixedPSDrop10\endcsname{3.13}
\expandafter\def\csname vXLCLFixedPSPlanMs10\endcsname{150.6}
\expandafter\def\csname vXLCLStagPSDrop10\endcsname{4.36}
\expandafter\def\csname vXLCLStagPSPlanMs10\endcsname{150.9}
\expandafter\def\csname vXLCLASUtil50\endcsname{90.7}
\expandafter\def\csname vXLCLFixedUtil50\endcsname{89.2}
\expandafter\def\csname vXLCLStagUtil50\endcsname{90.1}
\expandafter\def\csname vXLCLFixedMatchedUtil50\endcsname{90.1}
\expandafter\def\csname vXLCLRepairUtil50\endcsname{87.4}
\expandafter\def\csname vXLCLNoDelayUtil50\endcsname{89.5}
\expandafter\def\csname vXLCLFixedPSUtil50\endcsname{90.3}
\expandafter\def\csname vXLCLStagPSUtil50\endcsname{89.5}
\expandafter\def\csname vXLCLASDrop50\endcsname{10.27}
\expandafter\def\csname vXLCLASPlanMs50\endcsname{111.7}
\expandafter\def\csname vXLCLFixedDrop50\endcsname{20.41}
\expandafter\def\csname vXLCLFixedPlanMs50\endcsname{250.6}
\expandafter\def\csname vXLCLStagDrop50\endcsname{15.74}
\expandafter\def\csname vXLCLStagPlanMs50\endcsname{93.4}
\expandafter\def\csname vXLCLFixedMatchedDrop50\endcsname{16.28}
\expandafter\def\csname vXLCLFixedMatchedPlanMs50\endcsname{180.9}
\expandafter\def\csname vXLCLRepairDrop50\endcsname{8.43}
\expandafter\def\csname vXLCLRepairPlanMs50\endcsname{17.8}
\expandafter\def\csname vXLCLNoDelayDrop50\endcsname{18.98}
\expandafter\def\csname vXLCLNoDelayPlanMs50\endcsname{146.1}
\expandafter\def\csname vXLCLFixedPSDrop50\endcsname{13.78}
\expandafter\def\csname vXLCLFixedPSPlanMs50\endcsname{150.7}
\expandafter\def\csname vXLCLStagPSDrop50\endcsname{18.34}
\expandafter\def\csname vXLCLStagPSPlanMs50\endcsname{152.8}
\expandafter\def\csname vXLCLdUtilFixed\endcsname{1.54}
\expandafter\def\csname vXLCLdUtilFixedCI\endcsname{[1.20,\,1.88]}
\expandafter\def\csname vXLCLdUtilStag\endcsname{0.67}
\expandafter\def\csname vXLCLdUtilStagCI\endcsname{[0.40,\,0.95]}
\expandafter\def\csname vXLCLdUtilFixedMatched\endcsname{0.66}
\expandafter\def\csname vXLCLdUtilFixedMatchedCI\endcsname{[0.36,\,0.95]}
\expandafter\def\csname vXLCLdUtilNoDelay\endcsname{1.20}
\expandafter\def\csname vXLCLdUtilNoDelayCI\endcsname{[0.88,\,1.54]}
\expandafter\def\csname vXLCLdUtilFixedPS\endcsname{0.44}
\expandafter\def\csname vXLCLdUtilFixedPSCI\endcsname{[0.20,\,0.69]}
\expandafter\def\csname vXLCLdUtilStagPS\endcsname{1.20}
\expandafter\def\csname vXLCLdUtilStagPSCI\endcsname{[0.86,\,1.54]}
\expandafter\def\csname vXLNMissions\endcsname{90}
\expandafter\def\csname vXLNKappas\endcsname{3}
\expandafter\def\csname vXLEqRegretAS\endcsname{0.205}
\expandafter\def\csname vXLEqEffAS\endcsname{0.300}
\expandafter\def\csname vXLEqRegretBaseMin\endcsname{0.240}
\expandafter\def\csname vXLEqRegretBaseMax\endcsname{0.383}
\expandafter\def\csname vXLEqEffBaseMin\endcsname{0.462}
\expandafter\def\csname vXLEqEffBaseMax\endcsname{0.657}
\expandafter\def\csname vXLEqRegretRedMin\endcsname{15}
\expandafter\def\csname vXLEqRegretRedMax\endcsname{46}
\expandafter\def\csname vXLNeedFixedMs\endcsname{288}
\expandafter\def\csname vXLSpeedFixed\endcsname{1.9}
\expandafter\def\csname vXLNeedStagMs\endcsname{246}
\expandafter\def\csname vXLSpeedStag\endcsname{1.6}
\expandafter\def\csname vXLRedXlargeMin\endcsname{0}
\expandafter\def\csname vXLRedXlargeMax\endcsname{20}
\expandafter\def\csname vXLDropRedFixed\endcsname{50}
\expandafter\def\csname vXLDropRedStag\endcsname{35}
\expandafter\def\csname vXLDropRedNoDelay\endcsname{46}
\expandafter\def\csname vXLPlanRatio1\endcsname{1.8}
\expandafter\def\csname vXLCLUtilBaseMinK\endcsname{89.2}
\expandafter\def\csname vXLCLUtilBaseMaxK\endcsname{90.3}
\expandafter\def\csname vXLAllocASNone\endcsname{67.2}
\expandafter\def\csname vXLAllocASHigh\endcsname{330.9}
\expandafter\def\csname vXLAllocASShare\endcsname{16}
\expandafter\def\csname vXLAllocFixedPSNone\endcsname{150.0}
\expandafter\def\csname vXLAllocFixedPSHigh\endcsname{150.0}
\expandafter\def\csname vXLAllocFixedPSShare\endcsname{36}
\expandafter\def\csname vXLAllocStagPSNone\endcsname{103.3}
\expandafter\def\csname vXLAllocStagPSHigh\endcsname{180.4}
\expandafter\def\csname vXLAllocStagPSShare\endcsname{28}
\expandafter\def\csname vXLAllocNoneShare\endcsname{36}
\expandafter\def\csname vXLAUCLearned\endcsname{0.926}
\expandafter\def\csname vXLAUCOnline\endcsname{0.757}
\expandafter\def\csname vXLPosRate\endcsname{1.3}
\expandafter\def\csname vXLASMsXlarge\endcsname{123}
}{}
\IfFileExists{tables/static_numbers_scale.tex}{% copied from the static benchmark (conference pipeline); generated by experiments/13_scale_equal_compute.py -- do not edit
\expandafter\def\csname vScSmallMs\endcsname{18}
\expandafter\def\csname vScSmallOurs\endcsname{0.006}
\expandafter\def\csname vScSmallRedMin\endcsname{79}
\expandafter\def\csname vScSmallRedMax\endcsname{79}
\expandafter\def\csname vScSmallWinStatic\endcsname{100}
\expandafter\def\csname vScSmallSpeedMin\endcsname{2.7}
\expandafter\def\csname vScSmallSpeedMax\endcsname{2.8}
\expandafter\def\csname vScMediumMs\endcsname{25}
\expandafter\def\csname vScMediumOurs\endcsname{0.030}
\expandafter\def\csname vScMediumRedMin\endcsname{33}
\expandafter\def\csname vScMediumRedMax\endcsname{52}
\expandafter\def\csname vScMediumWinStatic\endcsname{92}
\expandafter\def\csname vScMediumSpeedMin\endcsname{1.5}
\expandafter\def\csname vScMediumSpeedMax\endcsname{2.0}
\expandafter\def\csname vScLargeMs\endcsname{81}
\expandafter\def\csname vScLargeOurs\endcsname{0.213}
\expandafter\def\csname vScLargeRedMin\endcsname{22}
\expandafter\def\csname vScLargeRedMax\endcsname{27}
\expandafter\def\csname vScLargeWinStatic\endcsname{63}
\expandafter\def\csname vScLargeSpeedMin\endcsname{1.7}
\expandafter\def\csname vScLargeSpeedMax\endcsname{2.1}
\expandafter\def\csname vScXlargeMs\endcsname{151}
\expandafter\def\csname vScXlargeOurs\endcsname{0.205}
\expandafter\def\csname vScXlargeRedMin\endcsname{12}
\expandafter\def\csname vScXlargeRedMax\endcsname{20}
\expandafter\def\csname vScXlargeWinStatic\endcsname{62}
\expandafter\def\csname vScXlargeSpeedMin\endcsname{1.7}
\expandafter\def\csname vScXlargeSpeedMax\endcsname{2.0}
\expandafter\def\csname vScAllMs\endcsname{42}
\expandafter\def\csname vScAllOurs\endcsname{0.083}
\expandafter\def\csname vScAllRedMin\endcsname{35}
\expandafter\def\csname vScAllRedMax\endcsname{41}
\expandafter\def\csname vScAllWinStatic\endcsname{91}
\expandafter\def\csname vScAllSpeedMin\endcsname{2.4}
\expandafter\def\csname vScAllSpeedMax\endcsname{3.4}
\expandafter\def\csname vClWinFixed\endcsname{84}
\expandafter\def\csname vClWinStag\endcsname{84}
\expandafter\def\csname vClWinFixedPS\endcsname{79}
}{}
\IfFileExists{tables/numbers.tex}{% generated by experiments/make_paper_numbers.py -- do not edit
\expandafter\def\csname vJCalibFlights\endcsname{18}
\expandafter\def\csname vJCalibRTF\endcsname{1.99}
\expandafter\def\csname vJDAS\endcsname{+0.28 [$-$0.22, +0.79]}
\expandafter\def\csname vJDASEvent\endcsname{+0.28 [$-$0.22, +0.79]}
\expandafter\def\csname vJDASFixedForty\endcsname{+0.12 [$-$0.32, +0.57]}
\expandafter\def\csname vJDFixed\endcsname{+0.34 [$-$0.30, +0.97]}
\expandafter\def\csname vJDFixedForty\endcsname{+0.40 [$-$0.06, +0.85]}
\expandafter\def\csname vJDFixedFortyEvent\endcsname{+0.40 [$-$0.06, +0.85]}
\expandafter\def\csname vJDMsAS\endcsname{$-$5.5 [$-$6.7, $-$4.2]}
\expandafter\def\csname vJDNoDelay\endcsname{+0.04 [$-$0.46, +0.56]}
\expandafter\def\csname vJDRepair\endcsname{+3.81 [+3.04, +4.60]}
\expandafter\def\csname vJDRiskZero\endcsname{+0.12 [$-$0.35, +0.60]}
\expandafter\def\csname vJDStag\endcsname{+0.27 [$-$0.23, +0.78]}
\expandafter\def\csname vJDmAS\endcsname{+0.28}
\expandafter\def\csname vJDmFixed\endcsname{+0.34}
\expandafter\def\csname vJDmFixedForty\endcsname{+0.40}
\expandafter\def\csname vJDmNoDelay\endcsname{+0.04}
\expandafter\def\csname vJDmRepair\endcsname{+3.81}
\expandafter\def\csname vJDmRiskZero\endcsname{+0.12}
\expandafter\def\csname vJDmStag\endcsname{+0.27}
\expandafter\def\csname vJLatC\endcsname{0.147}
\expandafter\def\csname vJLatCpn\endcsname{0.204}
\expandafter\def\csname vJLossAS\endcsname{0}
\expandafter\def\csname vJLossFixed\endcsname{0}
\expandafter\def\csname vJLossFixedForty\endcsname{0}
\expandafter\def\csname vJLossStag\endcsname{0}
\expandafter\def\csname vJMsAS\endcsname{55.3}
\expandafter\def\csname vJMsFixed\endcsname{250.2}
\expandafter\def\csname vJMsFixedForty\endcsname{40.2}
\expandafter\def\csname vJMsNoDelay\endcsname{61.8}
\expandafter\def\csname vJMsRepair\endcsname{1.1}
\expandafter\def\csname vJMsRisk\endcsname{49.9}
\expandafter\def\csname vJMsRiskZero\endcsname{56.4}
\expandafter\def\csname vJMsStag\endcsname{63.7}
\expandafter\def\csname vJNConfigsMain\endcsname{2}
\expandafter\def\csname vJNMain\endcsname{270}
\expandafter\def\csname vJNRunsMain\endcsname{2\,430}
\expandafter\def\csname vJPAS\endcsname{0.542}
\expandafter\def\csname vJPFixedForty\endcsname{0.268}
\expandafter\def\csname vJPStag\endcsname{0.542}
\expandafter\def\csname vJReproMatched\endcsname{24 of 24}
\expandafter\def\csname vJReproSpeedup\endcsname{21}
\expandafter\def\csname vJTrigMax\endcsname{+0.95}
\expandafter\def\csname vJTrigMin\endcsname{+0.95}
\expandafter\def\csname vJUAS\endcsname{77.76}
\expandafter\def\csname vJUFixed\endcsname{77.71}
\expandafter\def\csname vJUFixedForty\endcsname{77.64}
\expandafter\def\csname vJUNoDelay\endcsname{78.00}
\expandafter\def\csname vJURepair\endcsname{74.70}
\expandafter\def\csname vJURisk\endcsname{78.04}
\expandafter\def\csname vJURiskZero\endcsname{77.92}
\expandafter\def\csname vJUStag\endcsname{77.77}
\expandafter\def\csname vJValidatorPlans\endcsname{4\,779}
\expandafter\def\csname vJWinAS\endcsname{0}
\expandafter\def\csname vJWinFixed\endcsname{0}
\expandafter\def\csname vJWinFixedForty\endcsname{0}
\expandafter\def\csname vJWinStag\endcsname{0}
}{}
\IfFileExists{tables/numbers_s2.tex}{% generated by experiments/analyze_study2.py -- do not edit
\expandafter\def\csname vKBCTau\endcsname{1.00}
\expandafter\def\csname vKBCn\endcsname{2}
\expandafter\def\csname vKBCr\endcsname{1.00}
\expandafter\def\csname vKBgzDBuf\endcsname{+1.90 [$-$0.59, +4.65]}
\expandafter\def\csname vKBgzDConv\endcsname{+1.80 [$-$2.31, +5.65]}
\expandafter\def\csname vKBgzDForty\endcsname{+2.82 [$-$1.20, +6.62]}
\expandafter\def\csname vKCDForty\endcsname{+12.20 [+12.20, +12.20]}
\expandafter\def\csname vKCFails\endcsname{0}
\expandafter\def\csname vKCInvalid\endcsname{0}
\expandafter\def\csname vKCMsForty\endcsname{40.1}
\expandafter\def\csname vKCMsProp\endcsname{28.1}
\expandafter\def\csname vKCN\endcsname{1}
\expandafter\def\csname vKCPForty\endcsname{$<$0.001}
\expandafter\def\csname vKCRTF\endcsname{1.94}
\expandafter\def\csname vKCRTFmin\endcsname{1.93}
\expandafter\def\csname vKCRuns\endcsname{2}
\expandafter\def\csname vKCUForty\endcsname{48.78}
\expandafter\def\csname vKCUProp\endcsname{60.98}
\expandafter\def\csname vKCallsAS\endcsname{23.2}
\expandafter\def\csname vKCallsAblStop\endcsname{23.4}
\expandafter\def\csname vKCallsAblTrig\endcsname{6.8}
\expandafter\def\csname vKCallsAblZ\endcsname{23.3}
\expandafter\def\csname vKCallsBuf\endcsname{23.4}
\expandafter\def\csname vKCallsConv\endcsname{6.8}
\expandafter\def\csname vKCallsFix\endcsname{22.9}
\expandafter\def\csname vKCallsForty\endcsname{23.3}
\expandafter\def\csname vKCallsOnl\endcsname{23.2}
\expandafter\def\csname vKCallsProf\endcsname{23.2}
\expandafter\def\csname vKCallsProp\endcsname{23.3}
\expandafter\def\csname vKCallsRep\endcsname{23.2}
\expandafter\def\csname vKCallsRisk\endcsname{23.2}
\expandafter\def\csname vKCallsScaled\endcsname{23.3}
\expandafter\def\csname vKCallsStag\endcsname{23.2}
\expandafter\def\csname vKDAS\endcsname{+2.84 [+2.29, +3.40]}
\expandafter\def\csname vKDAblStop\endcsname{+0.22 [+0.07, +0.37]}
\expandafter\def\csname vKDAblStopKFifty\endcsname{+0.08 [$-$0.12, +0.29]}
\expandafter\def\csname vKDAblStopKTen\endcsname{+0.22 [+0.07, +0.38]}
\expandafter\def\csname vKDAblStopKTwoHundred\endcsname{+0.35 [+0.08, +0.64]}
\expandafter\def\csname vKDAblTrig\endcsname{+0.22 [+0.06, +0.38]}
\expandafter\def\csname vKDAblZ\endcsname{+0.18 [$-$0.06, +0.41]}
\expandafter\def\csname vKDAblZWHalfHigh\endcsname{$-$0.01 [$-$0.08, +0.07]}
\expandafter\def\csname vKDAblZWHalfMid\endcsname{$-$0.82 [$-$1.52, $-$0.11]}
\expandafter\def\csname vKDAblZWOneHigh\endcsname{+0.03 [$-$0.15, +0.21]}
\expandafter\def\csname vKDAblZWOneMid\endcsname{+1.51 [+0.86, +2.16]}
\expandafter\def\csname vKDBuf\endcsname{+1.35 [+0.93, +1.77]}
\expandafter\def\csname vKDBufKFifty\endcsname{+1.15 [+0.70, +1.61]}
\expandafter\def\csname vKDBufKTen\endcsname{+1.32 [+0.86, +1.79]}
\expandafter\def\csname vKDBufKTwoHundred\endcsname{+1.56 [+1.11, +2.01]}
\expandafter\def\csname vKDBufWHalfHigh\endcsname{$-$2.01 [$-$2.80, $-$1.21]}
\expandafter\def\csname vKDBufWHalfMid\endcsname{+2.48 [+1.98, +3.00]}
\expandafter\def\csname vKDBufWOneHigh\endcsname{+3.03 [+2.01, +4.11]}
\expandafter\def\csname vKDBufWOneMid\endcsname{+1.89 [+1.24, +2.52]}
\expandafter\def\csname vKDConv\endcsname{+4.44 [+3.82, +5.05]}
\expandafter\def\csname vKDConvKFifty\endcsname{+4.03 [+3.36, +4.68]}
\expandafter\def\csname vKDConvKTen\endcsname{+3.99 [+3.32, +4.64]}
\expandafter\def\csname vKDConvKTwoHundred\endcsname{+5.30 [+4.69, +5.91]}
\expandafter\def\csname vKDConvWHalfHigh\endcsname{+0.10 [$-$0.90, +1.07]}
\expandafter\def\csname vKDConvWHalfMid\endcsname{+3.34 [+2.42, +4.26]}
\expandafter\def\csname vKDConvWOneHigh\endcsname{+7.68 [+6.62, +8.73]}
\expandafter\def\csname vKDConvWOneMid\endcsname{+6.65 [+5.79, +7.51]}
\expandafter\def\csname vKDFix\endcsname{+5.92 [+5.41, +6.44]}
\expandafter\def\csname vKDFixKFifty\endcsname{+3.99 [+3.40, +4.58]}
\expandafter\def\csname vKDFixKTen\endcsname{+2.44 [+1.84, +3.03]}
\expandafter\def\csname vKDFixKTwoHundred\endcsname{+11.34 [+10.78, +11.91]}
\expandafter\def\csname vKDForty\endcsname{+2.87 [+2.29, +3.45]}
\expandafter\def\csname vKDFortyKFifty\endcsname{+2.71 [+2.10, +3.33]}
\expandafter\def\csname vKDFortyKTen\endcsname{+2.58 [+1.93, +3.22]}
\expandafter\def\csname vKDFortyKTwoHundred\endcsname{+3.33 [+2.74, +3.91]}
\expandafter\def\csname vKDFortyWHalfHigh\endcsname{$-$1.57 [$-$2.52, $-$0.64]}
\expandafter\def\csname vKDFortyWHalfMid\endcsname{+1.64 [+0.81, +2.45]}
\expandafter\def\csname vKDFortyWOneHigh\endcsname{+6.22 [+5.17, +7.29]}
\expandafter\def\csname vKDFortyWOneMid\endcsname{+5.20 [+4.36, +6.03]}
\expandafter\def\csname vKDOnl\endcsname{+2.70 [+2.13, +3.28]}
\expandafter\def\csname vKDOnlKFifty\endcsname{+2.77 [+2.15, +3.40]}
\expandafter\def\csname vKDOnlKTen\endcsname{+2.76 [+2.15, +3.36]}
\expandafter\def\csname vKDOnlKTwoHundred\endcsname{+2.57 [+1.97, +3.16]}
\expandafter\def\csname vKDProf\endcsname{+2.62 [+2.04, +3.19]}
\expandafter\def\csname vKDRep\endcsname{+4.21 [+3.57, +4.84]}
\expandafter\def\csname vKDRisk\endcsname{+2.49 [+1.92, +3.06]}
\expandafter\def\csname vKDScaled\endcsname{+2.57 [+1.99, +3.17]}
\expandafter\def\csname vKDScaledKFifty\endcsname{+2.57 [+1.95, +3.21]}
\expandafter\def\csname vKDScaledKTen\endcsname{+2.48 [+1.87, +3.10]}
\expandafter\def\csname vKDScaledKTwoHundred\endcsname{+2.66 [+2.04, +3.29]}
\expandafter\def\csname vKDStag\endcsname{+3.11 [+2.54, +3.69]}
\expandafter\def\csname vKDabsAS\endcsname{2.84}
\expandafter\def\csname vKDabsAblStop\endcsname{0.22}
\expandafter\def\csname vKDabsAblTrig\endcsname{0.22}
\expandafter\def\csname vKDabsAblZ\endcsname{0.18}
\expandafter\def\csname vKDabsBuf\endcsname{1.35}
\expandafter\def\csname vKDabsConv\endcsname{4.44}
\expandafter\def\csname vKDabsFix\endcsname{5.92}
\expandafter\def\csname vKDabsForty\endcsname{2.87}
\expandafter\def\csname vKDabsOnl\endcsname{2.70}
\expandafter\def\csname vKDabsProf\endcsname{2.62}
\expandafter\def\csname vKDabsRep\endcsname{4.21}
\expandafter\def\csname vKDabsRisk\endcsname{2.49}
\expandafter\def\csname vKDabsScaled\endcsname{2.57}
\expandafter\def\csname vKDabsStag\endcsname{3.11}
\expandafter\def\csname vKDmAS\endcsname{+2.84}
\expandafter\def\csname vKDmAblStop\endcsname{+0.22}
\expandafter\def\csname vKDmAblTrig\endcsname{+0.22}
\expandafter\def\csname vKDmAblZ\endcsname{+0.18}
\expandafter\def\csname vKDmBuf\endcsname{+1.35}
\expandafter\def\csname vKDmConv\endcsname{+4.44}
\expandafter\def\csname vKDmFix\endcsname{+5.92}
\expandafter\def\csname vKDmForty\endcsname{+2.87}
\expandafter\def\csname vKDmOnl\endcsname{+2.70}
\expandafter\def\csname vKDmProf\endcsname{+2.62}
\expandafter\def\csname vKDmRep\endcsname{+4.21}
\expandafter\def\csname vKDmRisk\endcsname{+2.49}
\expandafter\def\csname vKDmScaled\endcsname{+2.57}
\expandafter\def\csname vKDmStag\endcsname{+3.11}
\expandafter\def\csname vKHOneAll\endcsname{all}
\expandafter\def\csname vKHOneNSig\endcsname{6}
\expandafter\def\csname vKLateAS\endcsname{11.17}
\expandafter\def\csname vKLateAblStop\endcsname{5.74}
\expandafter\def\csname vKLateAblTrig\endcsname{5.85}
\expandafter\def\csname vKLateAblZ\endcsname{6.84}
\expandafter\def\csname vKLateBuf\endcsname{6.53}
\expandafter\def\csname vKLateBufWHalfHigh\endcsname{8.4}
\expandafter\def\csname vKLateBufWHalfMid\endcsname{3.6}
\expandafter\def\csname vKLateBufWOneHigh\endcsname{10.0}
\expandafter\def\csname vKLateBufWOneMid\endcsname{4.2}
\expandafter\def\csname vKLateConv\endcsname{10.31}
\expandafter\def\csname vKLateFix\endcsname{9.78}
\expandafter\def\csname vKLateForty\endcsname{11.22}
\expandafter\def\csname vKLateFortyWHalfHigh\endcsname{13.1}
\expandafter\def\csname vKLateFortyWHalfMid\endcsname{9.0}
\expandafter\def\csname vKLateFortyWOneHigh\endcsname{13.9}
\expandafter\def\csname vKLateFortyWOneMid\endcsname{8.9}
\expandafter\def\csname vKLateOnl\endcsname{11.17}
\expandafter\def\csname vKLateProf\endcsname{11.24}
\expandafter\def\csname vKLateProp\endcsname{5.72}
\expandafter\def\csname vKLatePropWHalfHigh\endcsname{6.2}
\expandafter\def\csname vKLatePropWHalfMid\endcsname{4.3}
\expandafter\def\csname vKLatePropWOneHigh\endcsname{8.1}
\expandafter\def\csname vKLatePropWOneMid\endcsname{4.3}
\expandafter\def\csname vKLateRedAS\endcsname{49}
\expandafter\def\csname vKLateRedAblStop\endcsname{0}
\expandafter\def\csname vKLateRedAblTrig\endcsname{2}
\expandafter\def\csname vKLateRedAblZ\endcsname{16}
\expandafter\def\csname vKLateRedBuf\endcsname{12}
\expandafter\def\csname vKLateRedConv\endcsname{44}
\expandafter\def\csname vKLateRedFix\endcsname{42}
\expandafter\def\csname vKLateRedForty\endcsname{49}
\expandafter\def\csname vKLateRedOnl\endcsname{49}
\expandafter\def\csname vKLateRedProf\endcsname{49}
\expandafter\def\csname vKLateRedRep\endcsname{49}
\expandafter\def\csname vKLateRedRisk\endcsname{49}
\expandafter\def\csname vKLateRedScaled\endcsname{49}
\expandafter\def\csname vKLateRedStag\endcsname{48}
\expandafter\def\csname vKLateRep\endcsname{11.19}
\expandafter\def\csname vKLateRisk\endcsname{11.21}
\expandafter\def\csname vKLateScaled\endcsname{11.18}
\expandafter\def\csname vKLateStag\endcsname{10.98}
\expandafter\def\csname vKLossBuf\endcsname{3}
\expandafter\def\csname vKLossConv\endcsname{0}
\expandafter\def\csname vKLossForty\endcsname{2}
\expandafter\def\csname vKLossOnl\endcsname{3}
\expandafter\def\csname vKLossProf\endcsname{3}
\expandafter\def\csname vKLossRisk\endcsname{3}
\expandafter\def\csname vKLossScaled\endcsname{3}
\expandafter\def\csname vKLostAS\endcsname{3.89}
\expandafter\def\csname vKLostAblStop\endcsname{0.00}
\expandafter\def\csname vKLostAblTrig\endcsname{0.00}
\expandafter\def\csname vKLostAblZ\endcsname{0.00}
\expandafter\def\csname vKLostBuf\endcsname{0.00}
\expandafter\def\csname vKLostConv\endcsname{3.32}
\expandafter\def\csname vKLostFix\endcsname{11.58}
\expandafter\def\csname vKLostFixKFifty\endcsname{7.96}
\expandafter\def\csname vKLostFixKTen\endcsname{3.34}
\expandafter\def\csname vKLostFixKTwoHundred\endcsname{23.42}
\expandafter\def\csname vKLostForty\endcsname{3.58}
\expandafter\def\csname vKLostFortyKFifty\endcsname{2.93}
\expandafter\def\csname vKLostFortyKTen\endcsname{2.10}
\expandafter\def\csname vKLostFortyKTwoHundred\endcsname{5.71}
\expandafter\def\csname vKLostOnl\endcsname{2.75}
\expandafter\def\csname vKLostProf\endcsname{3.52}
\expandafter\def\csname vKLostProp\endcsname{0.00}
\expandafter\def\csname vKLostPropKFifty\endcsname{0.00}
\expandafter\def\csname vKLostPropKTen\endcsname{0.00}
\expandafter\def\csname vKLostPropKTwoHundred\endcsname{0.00}
\expandafter\def\csname vKLostRep\endcsname{2.34}
\expandafter\def\csname vKLostRisk\endcsname{3.73}
\expandafter\def\csname vKLostScaled\endcsname{3.75}
\expandafter\def\csname vKLostStag\endcsname{5.50}
\expandafter\def\csname vKMsAS\endcsname{58.1}
\expandafter\def\csname vKMsAblStop\endcsname{40.1}
\expandafter\def\csname vKMsAblTrig\endcsname{57.8}
\expandafter\def\csname vKMsAblZ\endcsname{58.8}
\expandafter\def\csname vKMsBuf\endcsname{40.1}
\expandafter\def\csname vKMsConv\endcsname{40.2}
\expandafter\def\csname vKMsFix\endcsname{250.1}
\expandafter\def\csname vKMsForty\endcsname{40.1}
\expandafter\def\csname vKMsOnl\endcsname{11.2}
\expandafter\def\csname vKMsProf\endcsname{30.2}
\expandafter\def\csname vKMsProp\endcsname{61.0}
\expandafter\def\csname vKMsRep\endcsname{0.5}
\expandafter\def\csname vKMsRisk\endcsname{50.9}
\expandafter\def\csname vKMsScaled\endcsname{17.7}
\expandafter\def\csname vKMsStag\endcsname{55.8}
\expandafter\def\csname vKNConfigs\endcsname{12}
\expandafter\def\csname vKNMethods\endcsname{15}
\expandafter\def\csname vKNMissions\endcsname{270}
\expandafter\def\csname vKNRuns\endcsname{48\,600}
\expandafter\def\csname vKPBuf\endcsname{$<$0.001}
\expandafter\def\csname vKPForty\endcsname{$<$0.001}
\expandafter\def\csname vKPOnl\endcsname{$<$0.001}
\expandafter\def\csname vKPProf\endcsname{$<$0.001}
\expandafter\def\csname vKPRawAS\endcsname{$<$0.001}
\expandafter\def\csname vKPRawAblStop\endcsname{0.003}
\expandafter\def\csname vKPRawAblTrig\endcsname{0.006}
\expandafter\def\csname vKPRawAblZ\endcsname{0.142}
\expandafter\def\csname vKPRawConv\endcsname{$<$0.001}
\expandafter\def\csname vKPRawFix\endcsname{$<$0.001}
\expandafter\def\csname vKPRawRep\endcsname{$<$0.001}
\expandafter\def\csname vKPRawStag\endcsname{$<$0.001}
\expandafter\def\csname vKPRisk\endcsname{$<$0.001}
\expandafter\def\csname vKPScaled\endcsname{$<$0.001}
\expandafter\def\csname vKServedBufWHalfHigh\endcsname{25.3}
\expandafter\def\csname vKServedBufWHalfMid\endcsname{32.9}
\expandafter\def\csname vKServedBufWOneHigh\endcsname{36.7}
\expandafter\def\csname vKServedBufWOneMid\endcsname{49.1}
\expandafter\def\csname vKServedFortyWHalfHigh\endcsname{24.4}
\expandafter\def\csname vKServedFortyWHalfMid\endcsname{32.4}
\expandafter\def\csname vKServedFortyWOneHigh\endcsname{34.8}
\expandafter\def\csname vKServedFortyWOneMid\endcsname{47.0}
\expandafter\def\csname vKServedPropWHalfHigh\endcsname{24.3}
\expandafter\def\csname vKServedPropWHalfMid\endcsname{34.2}
\expandafter\def\csname vKServedPropWOneHigh\endcsname{39.0}
\expandafter\def\csname vKServedPropWOneMid\endcsname{50.4}
\expandafter\def\csname vKTostCI\endcsname{+0.22 [+0.09, +0.35]}
\expandafter\def\csname vKTostEquiv\endcsname{equivalent}
\expandafter\def\csname vKTostMs\endcsname{+20.9 [+17.5, +24.3]}
\expandafter\def\csname vKUAS\endcsname{49.31}
\expandafter\def\csname vKUAblStop\endcsname{51.93}
\expandafter\def\csname vKUAblTrig\endcsname{51.93}
\expandafter\def\csname vKUAblZ\endcsname{51.97}
\expandafter\def\csname vKUBuf\endcsname{50.81}
\expandafter\def\csname vKUConv\endcsname{47.71}
\expandafter\def\csname vKUFix\endcsname{46.23}
\expandafter\def\csname vKUForty\endcsname{49.28}
\expandafter\def\csname vKUOnl\endcsname{49.45}
\expandafter\def\csname vKUProf\endcsname{49.54}
\expandafter\def\csname vKUProp\endcsname{52.15}
\expandafter\def\csname vKUPropWHalfHigh\endcsname{33.38}
\expandafter\def\csname vKUPropWHalfMid\endcsname{47.65}
\expandafter\def\csname vKUPropWOneHigh\endcsname{55.94}
\expandafter\def\csname vKUPropWOneMid\endcsname{71.64}
\expandafter\def\csname vKURelAS\endcsname{5.8}
\expandafter\def\csname vKURelAblStop\endcsname{0.4}
\expandafter\def\csname vKURelAblTrig\endcsname{0.4}
\expandafter\def\csname vKURelAblZ\endcsname{0.3}
\expandafter\def\csname vKURelBuf\endcsname{2.6}
\expandafter\def\csname vKURelConv\endcsname{9.3}
\expandafter\def\csname vKURelFix\endcsname{12.8}
\expandafter\def\csname vKURelForty\endcsname{5.8}
\expandafter\def\csname vKURelOnl\endcsname{5.5}
\expandafter\def\csname vKURelProf\endcsname{5.3}
\expandafter\def\csname vKURelRep\endcsname{8.8}
\expandafter\def\csname vKURelRisk\endcsname{5.0}
\expandafter\def\csname vKURelScaled\endcsname{5.2}
\expandafter\def\csname vKURelStag\endcsname{6.3}
\expandafter\def\csname vKURep\endcsname{47.95}
\expandafter\def\csname vKURisk\endcsname{49.66}
\expandafter\def\csname vKUScaled\endcsname{49.58}
\expandafter\def\csname vKUSpreadStop\endcsname{0.38}
\expandafter\def\csname vKUStag\endcsname{49.04}
\expandafter\def\csname vKWinBuf\endcsname{9}
\expandafter\def\csname vKWinConv\endcsname{10}
\expandafter\def\csname vKWinForty\endcsname{8}
\expandafter\def\csname vKWinOnl\endcsname{9}
\expandafter\def\csname vKWinProf\endcsname{8}
\expandafter\def\csname vKWinRisk\endcsname{9}
\expandafter\def\csname vKWinScaled\endcsname{8}
}{}
\IfFileExists{tables/numbers_s3.tex}{% generated by experiments/analyze_study3.py -- do not edit
\expandafter\def\csname vLDBufAll\endcsname{+3.59 [+3.17, +4.02]}
\expandafter\def\csname vLDBufHighWHalf\endcsname{+3.31 [+2.58, +4.03]}
\expandafter\def\csname vLDBufHighWOne\endcsname{+3.91 [+2.83, +4.99]}
\expandafter\def\csname vLDBufMidWHalf\endcsname{+4.41 [+3.71, +5.13]}
\expandafter\def\csname vLDBufMidWOne\endcsname{+2.75 [+2.09, +3.39]}
\expandafter\def\csname vLDConvAll\endcsname{+6.19 [+5.69, +6.69]}
\expandafter\def\csname vLDConvHighWHalf\endcsname{+5.31 [+4.59, +6.02]}
\expandafter\def\csname vLDConvHighWOne\endcsname{+7.78 [+6.79, +8.77]}
\expandafter\def\csname vLDConvMidWHalf\endcsname{+5.43 [+4.68, +6.15]}
\expandafter\def\csname vLDConvMidWOne\endcsname{+6.25 [+5.48, +7.01]}
\expandafter\def\csname vLDFortyAll\endcsname{+4.61 [+4.15, +5.06]}
\expandafter\def\csname vLDFortyHighWHalf\endcsname{+3.61 [+2.94, +4.30]}
\expandafter\def\csname vLDFortyHighWOne\endcsname{+6.29 [+5.34, +7.23]}
\expandafter\def\csname vLDFortyMidWHalf\endcsname{+3.84 [+3.16, +4.51]}
\expandafter\def\csname vLDFortyMidWOne\endcsname{+4.70 [+3.95, +5.44]}
\expandafter\def\csname vLDMsProp\endcsname{$-$2.4 [$-$3.0, $-$1.9]}
\expandafter\def\csname vLDPropAll\endcsname{+1.97 [+1.61, +2.33]}
\expandafter\def\csname vLDPropBufHighWHalf\endcsname{$-$1.95 [$-$2.72, $-$1.20]}
\expandafter\def\csname vLDPropBufHighWOne\endcsname{+3.80 [+2.86, +4.72]}
\expandafter\def\csname vLDPropBufMidWHalf\endcsname{+2.40 [+1.85, +2.95]}
\expandafter\def\csname vLDPropBufMidWOne\endcsname{+2.23 [+1.59, +2.87]}
\expandafter\def\csname vLDPropHighWHalf\endcsname{+5.25 [+4.49, +6.05]}
\expandafter\def\csname vLDPropHighWOne\endcsname{+0.11 [$-$0.68, +0.89]}
\expandafter\def\csname vLDPropMidWHalf\endcsname{+2.01 [+1.32, +2.71]}
\expandafter\def\csname vLDPropMidWOne\endcsname{+0.52 [$-$0.01, +1.06]}
\expandafter\def\csname vLDPropRepBuf\endcsname{+1.62 [+1.22, +2.01]}
\expandafter\def\csname vLDPropRepForty\endcsname{+2.64 [+2.10, +3.15]}
\expandafter\def\csname vLDabsBufAll\endcsname{3.59}
\expandafter\def\csname vLDabsBufHighWHalf\endcsname{3.31}
\expandafter\def\csname vLDabsFortyAll\endcsname{4.61}
\expandafter\def\csname vLDabsPropAll\endcsname{1.97}
\expandafter\def\csname vLLateBuf\endcsname{6.62}
\expandafter\def\csname vLLateBufHighWHalf\endcsname{8.4}
\expandafter\def\csname vLLateBufHighWOne\endcsname{10.2}
\expandafter\def\csname vLLateBufMidWHalf\endcsname{3.5}
\expandafter\def\csname vLLateBufMidWOne\endcsname{4.3}
\expandafter\def\csname vLLateConv\endcsname{10.27}
\expandafter\def\csname vLLateConvHighWHalf\endcsname{12.0}
\expandafter\def\csname vLLateConvHighWOne\endcsname{12.7}
\expandafter\def\csname vLLateConvMidWHalf\endcsname{8.6}
\expandafter\def\csname vLLateConvMidWOne\endcsname{7.7}
\expandafter\def\csname vLLateFB\endcsname{8.33}
\expandafter\def\csname vLLateFBHighWHalf\endcsname{10.4}
\expandafter\def\csname vLLateFBHighWOne\endcsname{10.9}
\expandafter\def\csname vLLateFBMidWHalf\endcsname{6.2}
\expandafter\def\csname vLLateFBMidWOne\endcsname{5.9}
\expandafter\def\csname vLLateForty\endcsname{11.18}
\expandafter\def\csname vLLateFortyHighWHalf\endcsname{13.1}
\expandafter\def\csname vLLateFortyHighWOne\endcsname{14.0}
\expandafter\def\csname vLLateFortyMidWHalf\endcsname{9.1}
\expandafter\def\csname vLLateFortyMidWOne\endcsname{8.5}
\expandafter\def\csname vLLateProp\endcsname{5.70}
\expandafter\def\csname vLLatePropHighWHalf\endcsname{6.1}
\expandafter\def\csname vLLatePropHighWOne\endcsname{8.1}
\expandafter\def\csname vLLatePropMidWHalf\endcsname{4.2}
\expandafter\def\csname vLLatePropMidWOne\endcsname{4.3}
\expandafter\def\csname vLMsBuf\endcsname{40.1}
\expandafter\def\csname vLMsConv\endcsname{40.1}
\expandafter\def\csname vLMsFB\endcsname{55.0}
\expandafter\def\csname vLMsForty\endcsname{40.1}
\expandafter\def\csname vLMsProp\endcsname{57.4}
\expandafter\def\csname vLNConfigs\endcsname{12}
\expandafter\def\csname vLNMissions\endcsname{270}
\expandafter\def\csname vLNRuns\endcsname{16\,200}
\expandafter\def\csname vLNSig\endcsname{4}
\expandafter\def\csname vLPBufAll\endcsname{$<$0.001}
\expandafter\def\csname vLPBufHighWHalf\endcsname{$<$0.001}
\expandafter\def\csname vLPFortyAll\endcsname{$<$0.001}
\expandafter\def\csname vLPPropAll\endcsname{$<$0.001}
\expandafter\def\csname vLServedBufHighWHalf\endcsname{25.7}
\expandafter\def\csname vLServedBufHighWOne\endcsname{37.0}
\expandafter\def\csname vLServedBufMidWHalf\endcsname{33.3}
\expandafter\def\csname vLServedBufMidWOne\endcsname{49.1}
\expandafter\def\csname vLServedConvHighWHalf\endcsname{23.8}
\expandafter\def\csname vLServedConvHighWOne\endcsname{34.3}
\expandafter\def\csname vLServedConvMidWHalf\endcsname{31.6}
\expandafter\def\csname vLServedConvMidWOne\endcsname{46.3}
\expandafter\def\csname vLServedFBHighWHalf\endcsname{26.6}
\expandafter\def\csname vLServedFBHighWOne\endcsname{38.7}
\expandafter\def\csname vLServedFBMidWHalf\endcsname{35.0}
\expandafter\def\csname vLServedFBMidWOne\endcsname{50.5}
\expandafter\def\csname vLServedFortyHighWHalf\endcsname{25.0}
\expandafter\def\csname vLServedFortyHighWOne\endcsname{35.4}
\expandafter\def\csname vLServedFortyMidWHalf\endcsname{32.9}
\expandafter\def\csname vLServedFortyMidWOne\endcsname{47.7}
\expandafter\def\csname vLServedPropHighWHalf\endcsname{25.0}
\expandafter\def\csname vLServedPropHighWOne\endcsname{39.3}
\expandafter\def\csname vLServedPropMidWHalf\endcsname{34.8}
\expandafter\def\csname vLServedPropMidWOne\endcsname{50.8}
\expandafter\def\csname vLUBuf\endcsname{50.80}
\expandafter\def\csname vLUBufHighWHalf\endcsname{35.92}
\expandafter\def\csname vLUBufHighWOne\endcsname{52.48}
\expandafter\def\csname vLUBufMidWHalf\endcsname{45.48}
\expandafter\def\csname vLUBufMidWOne\endcsname{69.33}
\expandafter\def\csname vLUConv\endcsname{48.20}
\expandafter\def\csname vLUConvHighWHalf\endcsname{33.91}
\expandafter\def\csname vLUConvHighWOne\endcsname{48.61}
\expandafter\def\csname vLUConvMidWHalf\endcsname{44.46}
\expandafter\def\csname vLUConvMidWOne\endcsname{65.83}
\expandafter\def\csname vLUFB\endcsname{54.40}
\expandafter\def\csname vLUFBHighWHalf\endcsname{39.22}
\expandafter\def\csname vLUFBHighWOne\endcsname{56.39}
\expandafter\def\csname vLUFBMidWHalf\endcsname{49.89}
\expandafter\def\csname vLUFBMidWOne\endcsname{72.08}
\expandafter\def\csname vLUForty\endcsname{49.79}
\expandafter\def\csname vLUFortyHighWHalf\endcsname{35.61}
\expandafter\def\csname vLUFortyHighWOne\endcsname{50.10}
\expandafter\def\csname vLUFortyMidWHalf\endcsname{46.05}
\expandafter\def\csname vLUFortyMidWOne\endcsname{67.38}
\expandafter\def\csname vLUProp\endcsname{52.42}
\expandafter\def\csname vLUPropHighWHalf\endcsname{33.97}
\expandafter\def\csname vLUPropHighWOne\endcsname{56.28}
\expandafter\def\csname vLUPropMidWHalf\endcsname{47.88}
\expandafter\def\csname vLUPropMidWOne\endcsname{71.56}
\expandafter\def\csname vLURelConv\endcsname{12.8}
}{}
\IfFileExists{tables/numbers_s3x.tex}{% generated by experiments/paper_study3_extras.py -- do not edit
\expandafter\def\csname vLLateIncProp\endcsname{46}
\expandafter\def\csname vLLateRedForty\endcsname{25}
\expandafter\def\csname vLLostBuf\endcsname{0.00}
\expandafter\def\csname vLLostConv\endcsname{3.00}
\expandafter\def\csname vLLostFB\endcsname{1.55}
\expandafter\def\csname vLLostForty\endcsname{3.29}
\expandafter\def\csname vLLostProp\endcsname{0.00}
\expandafter\def\csname vLServedBuf\endcsname{36.3}
\expandafter\def\csname vLServedConv\endcsname{34.0}
\expandafter\def\csname vLServedFB\endcsname{37.7}
\expandafter\def\csname vLServedForty\endcsname{35.3}
\expandafter\def\csname vLServedProp\endcsname{37.5}
\expandafter\def\csname vLSkipBuf\endcsname{13.61}
\expandafter\def\csname vLSkipConv\endcsname{15.41}
\expandafter\def\csname vLSkipFB\endcsname{12.51}
\expandafter\def\csname vLSkipForty\endcsname{37.38}
\expandafter\def\csname vLSkipProp\endcsname{6.70}
\expandafter\def\csname vLURelBuf\endcsname{7.1}
\expandafter\def\csname vLURelConv\endcsname{12.8}
\expandafter\def\csname vLURelForty\endcsname{9.3}
\expandafter\def\csname vLURelProp\endcsname{3.8}
}{}
\newcommand{\val}[1]{\ifcsname v#1\endcsname\csname v#1\endcsname\else\todo{#1}\fi}
\newcommand{\fig}[2]{\IfFileExists{figures/#1.pdf}{\includegraphics[width=#2]{figures/#1.pdf}}{\fbox{\todo{figure \detokenize{#1}}}}}
\newcommand{\inputtable}[1]{\IfFileExists{tables/#1.tex}{\input{tables/#1.tex}}{\todo{table \detokenize{#1}}}}
\renewcommand{\thesection}{S\arabic{section}}
\renewcommand{\thetable}{S\arabic{table}}
\renewcommand{\thefigure}{S\arabic{figure}}
\renewcommand{\theequation}{S\arabic{equation}}
\journal{Robotics and Autonomous Systems}

\begin{document}
\begin{frontmatter}
\title{Supplementary material for: Closed-form chance constraints for fast multi-UAV replanning under uncertain
execution}
\author{Cao-Phuc Ha, Anh Nguyen Vu Minh, Phu-Nguyen Le}
\begin{abstract}
This document contains the parts of the evaluation that the main paper only summarises: the stopping rule and its two
delay costs with their properties, the scenario generator and the fitted error models, the exact-execution benchmark of
the stopping rules, the full results of study~1, and further results of study~2 per regime and in tier C. Section,
equation and table numbers without the prefix S refer to the main paper.
\end{abstract}
\end{frontmatter}

\section{Stopping rule and delay costs}\label{sec:Sstop}
\subsection{Learned stopping rule}
Every 5\,ms, the stopping rule compares the predicted gain of searching for another $H=10$\,ms with the cost of doing so.
With $\widehat{\Delta J}$ as the predicted gain, the search stops at time $\tau$ if
\begin{equation}\label{eq:Sstop}
\widehat{\Delta J}\;\le\;C_{\mathrm{delay}}(\tau)+\mu W H .
\end{equation}
The gain is predicted from the search history by a Poisson generalised linear model with a log
link~\cite{mccullagh1989glm} and a small $L_2$ penalty. Its eleven features describe the elapsed time, the number of
iterations, the recent improvements over 5\,ms, over 20\,ms and over the second half of the elapsed time, the time since
the last improvement, the total improvement over the repaired plan, the unassigned share of the reward, the numbers of
tasks and UAVs, and the iteration rate. The target is the improvement in the next 10\,ms in percent of $W$. The model is
fitted on 500\,ms search trajectories of development scenarios only, and a prediction costs one dot product. The search
loop enforces the hard time budget $B$ independently of the stopping rule. The price of computation is
$\mu=10^{-5}$ per millisecond.

\subsection{Deterministic delay cost}
The delay cost is computed from the best plan found so far. A late start of a route is partly absorbed by the waiting
times along it, in the same way as the forward time slack of routing with time windows~\cite{savelsbergh1992vrptw}. The
tolerance $\mathrm{tol}_j$ of a task is the margin between its finish time and its deadline, plus all waiting times on
the route up to this task. A task is lost when the start shift $\sigma_j(\tau)$ of its route exceeds its tolerance. The
start shift is zero for a UAV that is still busy, and it equals the waiting time until the publication for an idle UAV.
With $v_j=\alpha w_j-\delta\,\mathbb{1}[\text{$j$ changes UAV}]$ as the value of task $j$ in the objective and
$\Lambda(\tau)=\sum_j v_j\,\mathbb{1}[\sigma_j(\tau)>\mathrm{tol}_j]$, the delay cost of searching $H$ milliseconds longer
is $C_{\mathrm{delay}}(\tau)=\Lambda(\tau+H)-\Lambda(\tau)$.

\begin{prop}\label{prop:delay}
For a feasible plan executed exactly, (i) $\Lambda(\tau)$ is the value lost through missed deadlines if every UAV
follows its route, (ii) the loss of the objective is at most $\Lambda(\tau)$ if late tasks are skipped, and
(iii) $\Lambda(\tau)$ does not decrease with $\tau$ and is zero while all UAVs are busy with their committed tasks.
Hence $C_{\mathrm{delay}}$ is non-negative and vanishes while no UAV is waiting.
\end{prop}
\begin{proof}
A route that starts later by some shift reaches its first task later by the same amount. Each waiting time absorbs part
of the delay, so the delay at the end of a task equals the start shift minus all earlier waiting times, if this
difference is positive. The task therefore misses its deadline exactly when the start shift exceeds its tolerance.
Waiting consumes no energy, which proves (i). If late tasks are skipped, the triangle inequality ensures that every
later task is reached no later and with no more energy. All tasks that are on time in (i) remain on time, and the flown
distance, hence the energy term $\gamma E$, does not grow. The objective therefore loses the values $v_j$ of the skipped
tasks and nothing else, which is at most $\Lambda(\tau)$ and proves (ii). Part (iii) holds because the start shift does
not decrease with $\tau$ and every tolerance is non-negative in a feasible plan.
\end{proof}

\subsection{Execution-aware delay cost}\label{sec:Srisk}
Under uncertain execution, task $j$ is lost when $\sigma_j(\tau)+e_j>\mathrm{tol}_j$, where $e_j$ is its finish-time
error. The execution-aware cost replaces the indicator in $\Lambda$ by the probability of this event. The error of task
$j$ is modelled as $e_j\sim\mathcal{N}(m(h_j),s(h_j)^2)$, where $h_j=f_j-t_0$ is its lead time, the nominal finish time
minus the time of the planning snapshot. The mean and the spread are linear in the lead time,
\begin{equation}\label{eq:Serrmodel}
m(h)=b_0+b_1h,\qquad s(h)=s_0+s_1h,
\end{equation}
because errors accumulate along a route. The expected dissemination latency $\bar L$ is added to the start shift. The
expected loss and the execution-aware delay cost are
\begin{equation}\label{eq:Srisk}
\Lambda_{\mathrm{R}}(\tau)=\sum_j v_j\,\Phi\!\left(\frac{\sigma_j(\tau)+\bar L+m(h_j)-\mathrm{tol}_j}{s(h_j)}\right),
\qquad C_{\mathrm{R}}(\tau)=\Lambda_{\mathrm{R}}(\tau+H)-\Lambda_{\mathrm{R}}(\tau),
\end{equation}
where $\Phi$ is the standard normal distribution function. Search and predictor are unchanged, and $C_{\mathrm{R}}$
replaces $C_{\mathrm{delay}}$ in Eq.~\eqref{eq:Sstop}.

\begin{prop}\label{prop:risk}
(i) With $b_0=b_1=s_0=s_1=0$ and $\bar L=0$, $\Lambda_{\mathrm{R}}=\Lambda$, and the stopping rule takes exactly the
decisions of the deterministic one. (ii) $\Lambda_{\mathrm{R}}(\tau)$ does not decrease with $\tau$, so $C_{\mathrm{R}}\ge0$.
(iii) Under the error model, $\Lambda_{\mathrm{R}}(\tau)$ is the expected value of the tasks whose deadline is missed if
every UAV follows its route.
\end{prop}
\begin{proof}
With zero spread, $\Phi$ is read as the indicator of a positive argument, which gives (i). The start shift does not
decrease with $\tau$ and $\Phi$ is non-decreasing, which gives (ii). Part (iii) follows from the proof of
Proposition~\ref{prop:delay}, because task $j$ misses its deadline exactly when its delay at the finish,
$\sigma_j+e_j$ minus the absorbed waiting time, is positive, and the expectation of the indicator of this event is the
probability in Eq.~\eqref{eq:Srisk}.
\end{proof}

The parameters of Eq.~\eqref{eq:Serrmodel} are fitted on development missions. The mean is estimated by least squares,
and the spread by least squares on the absolute residuals scaled by $\sqrt{\pi/2}$, which is less sensitive to the heavy
tail of the errors than squared residuals. The latency $\bar L$ is the mean of the latency distribution in tier B and
the measured mean plan-acceptance latency in tier C.



\section{Scenario generator, error models and exact-execution snapshots}\label{sec:Sscen}
\subsection{Scenarios}
Missions are created by a seeded generator with the parameters in Table~\ref{tab:params}. The map is a square with a
side of 10\,km and the depot at its centre, and half of the tasks are clustered around three hotspots. Half of the tasks
are known at the start, and the others arrive in batches. Three fleet scales with 5, 10 and 20 UAVs and 25, 50 and 140
tasks are combined with three event profiles. The arrivals profile contains six arrival batches, the failures profile
three batches and two UAV failures, and the deadlines profile three batches and four tightened deadlines. A fourth scale
with 40 UAVs and 280 tasks is used for scalability only.

\begin{table}[t]
\caption{Scenario generator parameters (distances in m, times in s, energy in metres of flight).}\label{tab:params}
\centering\small
\setlength{\tabcolsep}{4pt}
\begin{tabular}{ll@{\qquad}ll}
\toprule
Map side / horizon & 10\,000 / 1\,800 & Service time & $U[30,90]$ \\
UAV speed & $U[10,16]$\,m/s & Time-window slack & $U[180,600]$ \\
Energy capacity & $U[18,28]$\,km & Release delay & $U[0,120]$ \\
Energy reserve & 10\% of capacity & Rewards 1, 2, 3, 5 & probabilities .4, .3, .2, .1 \\
Hover cost & 6 per second & Hotspots & 3, std.\ 600, 50\% of tasks \\
Tasks known at start & 50\% & Scales (UAVs/tasks) & 5/25, 10/50, 20/140, 40/280 \\
\bottomrule
\end{tabular}
\end{table}

\subsection{Static snapshots with exact execution}
The stopping rules are first compared without execution uncertainty. Each mission is simulated once with a
deterministic reference planner, and every replanning call after the start is stored with its previous plan. All
methods are evaluated on these snapshots, and regret is measured against the best known objective of each snapshot,
which includes a 3\,s ALNS run and, for the small and medium scales, the OR-Tools routing solver~\cite{ortools} with
guided local search~\cite{voudouris1999gls}. Because all ALNS variants follow the same search trajectory, whole families
of each rule are evaluated exactly offline, with 12 fixed budgets, 10 stagnation patience values and 9 compute prices,
and each baseline is compared with AnytimeSwarm at the same mean computation. This benchmark uses its own pre-registered
test seeds 2000--2029 (\val{NsnapTest} snapshots of 270 missions), and it was reproduced on the final software platform
before the execution-aware study (Section~\ref{sec:repro}).



\inputtable{tab_errmodels}

\section{Exact execution: stopping rules at equal computation}\label{sec:Sstatic}
\label{sec:static}
On the \val{NsnapTest} test snapshots, regret-2 repair finished within about one millisecond but left a regret of
\val{RepairRegret}\%. With a budget of 250\,ms, ALNS with the full budget reached \val{FixedRegret}\%, whereas
AnytimeSwarm reached \val{ASRegret}\% with \val{ASUsed}\% of the budget, that is, with \val{SavedPct}\% less computation.
Because the stopping rules use different amounts of computation, they are best compared at equal computation.
Table~\ref{tab:scale} makes this comparison for each fleet scale, with the fixed budget and the stagnation rule evaluated
at the mean computation that AnytimeSwarm used on the same scale. AnytimeSwarm reached a lower regret on every scale,
and all paired intervals lie below zero. The reduction was \val{ScMediumRedMin}\% to \val{ScMediumRedMax}\% on the
medium scale, \val{ScLargeRedMin}\% to \val{ScLargeRedMax}\% on the large scale and \val{ScXlargeRedMin}\% to
\val{ScXlargeRedMax}\% with 40 UAVs, a scale that was never used for training. Read in the other direction, the two
baselines needed between \val{ScMediumSpeedMin} and \val{ScSmallSpeedMax} times the computation of AnytimeSwarm to
reach its regret. Over all scales, AnytimeSwarm reached \val{EqRegretAS}\% against \val{EqRegretBaseMin}\% to
\val{EqRegretBaseMax}\% for four baselines that include a fixed budget and a stagnation patience tuned per fleet scale
on development data.

\inputtable{tab_static_scale}

Fig.~\ref{fig:static} shows where this advantage comes from. On the medium and large scales, the quality--compute
frontier of AnytimeSwarm lies below those of fixed budgets and stagnation stopping over the range of practical
computation (Fig.~\ref{fig:static}a). Grouped by the improvement that a 500\,ms search can still find
(Fig.~\ref{fig:static}b), AnytimeSwarm spent \val{AllocASNone}\,ms on snapshots that cannot be improved and
\val{AllocASHigh}\,ms on those with the largest available improvement. Only \val{AllocASShare}\% of its computation went
to snapshots without any possible gain, compared with more than half for the baselines. The learned predictor is
essential for this allocation. It reached an area under the ROC curve~\cite{fawcett2006roc} of \val{AUCLearned} for the
event that the next 10\,ms bring an improvement, compared with \val{AUCOnline} for an extrapolation of the recent
improvement rate.

\begin{figure}[t]
\centering
\begin{minipage}[b]{0.62\linewidth}\centering\fig{fig_static_frontier}{\linewidth}\\[-1pt]{\small (a)}\end{minipage}\hfill
\begin{minipage}[b]{0.36\linewidth}\centering\fig{fig_static_allocation}{\linewidth}\\[-1pt]{\small (b)}\end{minipage}
\caption{Exact execution, test seeds 2000--2029: (a) quality--compute frontiers on identical search trajectories (open
markers: baselines tuned per scale; star: selected AnytimeSwarm setting); (b) mean computation per snapshot, grouped by
the improvement that a 500\,ms search can still find.}\label{fig:static}
\end{figure}

\subsection{Reproduction on the evaluation platform}\label{sec:repro}
The exact-execution benchmark was produced on Windows. Before the execution-aware study, its development pipeline was
re-run on the platform of all later experiments (Ubuntu 24.04 under WSL2, identical Python and package versions). The
deterministic part was reproduced bit for bit, except for differences of one unit in the last place of some energy
values, which stem from the floating-point libraries of the two operating systems. The search ran
\val{JReproSpeedup}\% more iterations per millisecond on the new platform. The improvement predictor, re-trained on the
new timing traces, reached the same compute price by the pre-registered rule, and the conclusions of the
equal-computation comparison agreed in sign and significance in \val{JReproMatched} comparisons. The frozen predictor
was used unchanged in all later experiments.


\section{Study 1: execution-aware stopping on the nominal planning model}\label{sec:Sstudy1}
Study~1 compares AnytimeSwarm-R with the stopping rules of the first group of Section~\ref{sec:methods} on the nominal
planning model, on held-out seeds 3000--3029 in tier B (270 missions over three scales and three profiles, and 90
missions with 40 UAVs) and 4000--4009 in tier C (30 missions per method). Its primary matrix crosses two execution
levels (low, mid), two latency levels (none, high), two latency factors ($\kappa_B=10, 50$) and two triggers (event,
periodic every 60\,s). Its primary hypothesis is that AnytimeSwarm-R serves more reward than ALNS with 40\,ms,
AnytimeSwarm and the stagnation rule.
\subsection{Tier B primary matrix}\label{sec:tierb}
The primary matrix comprises \val{JNRunsMain} mission runs, that is, \val{JNMain} held-out missions in
\val{JNConfigsMain} configurations with eight methods. Table~\ref{tab:main} lists utility, lost tasks and computation per
method, and Fig.~\ref{fig:main} relates utility to computation and shows the difference in every configuration.

Search pays off clearly under uncertain execution. Regret-2 repair reached a utility of \val{JURepair}\%, and
AnytimeSwarm-R reached \val{JURisk}\%, which is \val{JDmRepair} percentage points (pp) or \val{JURelRepair}\% more. The
full budget of 250\,ms was the weakest search method, with \val{JUFixed}\%, and AnytimeSwarm-R served \val{JDFixed}\,pp
more reward than it. The loss of the long budget stems from late plans. A task counts as lost to delay when it is part of
the computed plan but can no longer be served on time when the plan is released. With $\kappa_B=50$, the full budget
lost \val{JLostKFiftyFixed} tasks per mission in this way, against \val{JLostKFiftyRisk} for AnytimeSwarm-R, that is,
\val{JLostRedFixed}\% fewer. Without its delay cost, AnytimeSwarm lost \val{JLostKFiftyNoDelay} tasks per mission in this
way, but its utility was not lower (\val{JUNoDelay}\% against \val{JUAS}\%). The delay cost thus prevents losses from
late plans without a measurable gain in utility on this matrix.

Among the remaining search methods, the differences are small. All seven search methods lie within
\val{JUSpreadSearch}\,pp of each other. AnytimeSwarm-R reached the highest mean utility with the least computation of the adaptive rules,
\val{JMsRisk}\,ms per call against \val{JMsAS}\,ms for AnytimeSwarm and \val{JMsStag}\,ms for the stagnation rule. The
reduction of \val{JMsRedAS}\% relative to AnytimeSwarm is significant (\val{JDMsAS}\,ms per call), so the second
hypothesis holds. The primary hypothesis does not hold at the pre-registered level. The pooled advantage of
AnytimeSwarm-R was \val{JDFixedForty}\,pp over ALNS with 40\,ms, \val{JDAS}\,pp over AnytimeSwarm and \val{JDStag}\,pp
over the stagnation rule, with Holm-adjusted $p$-values of \val{JPFixedForty}, \val{JPAS} and \val{JPStag}. Only the
first interval excludes zero before correction. Per configuration and without correction for multiple testing,
AnytimeSwarm-R was better than ALNS with 40\,ms in \val{JWinFixedForty} of the \val{JNConfigsMain} configurations, better
than the stagnation rule in \val{JWinStag} and never different from AnytimeSwarm, which is consistent with chance
(Fig.~\ref{fig:main}b). The ablation that keeps only the latency term reached \val{JURiskZero}\%, close to AnytimeSwarm.
Together with the misspecification runs below, this suggests that a non-zero error spread mainly makes the rule stop
earlier, rather than that the calibration of the spread matters.

The fixed budget of 40\,ms is a strong reference, because it was matched to the mean computation of AnytimeSwarm on
development data. Its weakness lies in the allocation across fleet scales. AnytimeSwarm-R used \val{JMsScSmallRisk}\,ms
per call on the small scale and \val{JMsScLargeRisk}\,ms on the large scale. On the large scale, it served
\val{JDScLargeFixedForty}\,pp more reward than the budget of 40\,ms and was on a par with the full budget
(\val{JDScLargeFixed}\,pp). On the small scale, it served \val{JDScSmallFixed}\,pp more than the full budget and was on a
par with the budget of 40\,ms (\val{JDScSmallFixedForty}\,pp). A single fixed budget is therefore either too long for
small fleets or too short for large ones, whereas the adaptive rules find a suitable amount of computation on each scale
without tuning. The advantage over ALNS with 40\,ms was \val{JDFFKappa10}\,pp for $\kappa_B=10$ and
\val{JDFFKappa50}\,pp for $\kappa_B=50$, and \val{JDFFExecLow}\,pp under low and \val{JDFFExecMid}\,pp under mid
disturbances.

\inputtable{tab_main}

\begin{figure}[t]
\centering
\fig{fig_main}{\linewidth}
\caption{Tier B primary matrix, held-out seeds 3000--3029: (a) mean utility against mean planning time per call over
all \val{JNConfigsMain} configurations (regret-2 repair, at \val{JURepair}\% and 1\,ms, is not shown); (b) paired
difference of AnytimeSwarm-R minus ALNS with 40\,ms and minus AnytimeSwarm per configuration (execution level, latency
level, $\kappa_B$, trigger) with 95\% bootstrap intervals over \val{JNMain} missions.}\label{fig:main}
\end{figure}

\subsection{Replanning triggers}\label{sec:res-triggers}
The trigger mattered more than the stopping rule. Replanning every 60\,s in addition to the events raised the number of
planning calls from \val{JCallsEvent} to \val{JCallsPeriodic} per mission and the utility of every method, by
\val{JTrigMin}\,pp to \val{JTrigMax}\,pp, with all intervals above zero (\val{JTrigRisk}\,pp for AnytimeSwarm-R,
\val{JTrigFixedForty}\,pp for ALNS with 40\,ms). The third hypothesis therefore holds. Periodic replanning probably helps
because it corrects the accumulated deviations of the executed routes, which the event trigger only notices at the next
event. The advantage of AnytimeSwarm-R over ALNS with 40\,ms was \val{JDFixedFortyEvent}\,pp with the event trigger and
\val{JDFixedFortyPeriodic}\,pp with the periodic trigger. The two values are close, and the protocol did not test their
difference.

\subsection{Scalability to 40 UAVs}\label{sec:res-xl}
With 40 UAVs and 280 tasks, \val{JNXL} held-out missions were run under mid disturbances in \val{JXLNConf}
configurations of latency level and latency factor. All five methods reached between \val{JXLUAS}\% and
\val{JXLUStag}\%. AnytimeSwarm-R reached \val{JXLURisk}\% with \val{JXLMsRisk}\,ms per call. Its pooled difference was
\val{JXLDFixedOneEighty}\,pp to ALNS with 180\,ms, \val{JXLDAS}\,pp to AnytimeSwarm and \val{JXLDStag}\,pp to the
stagnation rule. After Holm correction over the four comparisons of the pre-registered analysis, none of them is
significant ($p_{\mathrm{Holm}}=$ \val{JXLPAS} for AnytimeSwarm), so the fifth hypothesis is not supported either.
AnytimeSwarm-R again needed less computation than AnytimeSwarm (\val{JXLDMsAS}\,ms per call). The stagnation rule
stopped earliest, after \val{JXLMsStag}\,ms, and was significantly better than AnytimeSwarm-R in \val{JXLLossStag} of the
\val{JXLNConf} configurations, which had no latency and $\kappa_B=50$. On this scale, the rules used between \val{JXLMsStag} and 250\,ms
per call, and utility hardly changed within this range of computation. The equal-computation advantage that
AnytimeSwarm showed with 40 UAVs under exact execution (Table~\ref{tab:scale}) therefore did not translate into a
utility gain under uncertain execution.

\subsection{Sensitivity analyses}\label{sec:res-sens}
The sensitivity runs use the held-out missions of the primary matrix without latency and are descriptive. They vary the
execution level beyond the primary matrix, the latency factor and the trigger in \val{JNSensConf} configurations, and
the error model in \val{JNMisConf} further configurations.

Under high disturbances, which lie outside the range of the primary matrix but within that of the fitted error models,
utility fell to about \val{JSensUHighRisk}\% for all methods. The combination of high disturbances and a slow planner
($\kappa_B=50$) gave the largest advantage of the risk-aware cost in this study. AnytimeSwarm-R served
\val{JSensHighKFiftyAS}\,pp more reward than AnytimeSwarm and used \val{JSensMsHighKFiftyRisk}\,ms instead of
\val{JSensMsHighKFiftyAS}\,ms per call. Its difference to ALNS with 40\,ms was \val{JSensHighKFiftyFixedForty}\,pp and to
the stagnation rule \val{JSensHighKFiftyStag}\,pp, so the gain over AnytimeSwarm brought it to the level of the
baselines rather than beyond them. With $\kappa_B=1$, planning time hardly delays the plan, both delay costs are close to
zero because they scale with the shift $\kappa\tau$, and the two rules behave alike. AnytimeSwarm-R then differed from
AnytimeSwarm by \val{JSensKOneAS}\,pp and \val{JSensKOneMsDiff}\,ms per call.

Replanning when a UAV deviated by more than 30\,s from its planned finish time raised the number of calls from
\val{JSensCallsEvent} to \val{JSensCallsDev} per mission and the utility by \val{JSensTrigDevRisk}\,pp for
AnytimeSwarm-R and \val{JSensTrigDevFixedForty}\,pp for ALNS with 40\,ms. Like the periodic trigger, it helped both
stopping rules by a similar amount. Under the deviation trigger, the difference of AnytimeSwarm-R was
\val{JSensDevFixedForty}\,pp to ALNS with 40\,ms and \val{JSensDevAS}\,pp to AnytimeSwarm.

The utility of AnytimeSwarm-R was insensitive to the error model. With the error model of the low, mid or high level
under low and mid disturbances, its difference to AnytimeSwarm ranged between \val{JMisAllLow}\,pp and
\val{JMisAllMid}\,pp, the correctly specified model was not better than a wrong one, and the mean utilities of all
versions lay within \val{JMisSpread}\,pp. Setting the error model to zero gave
\val{JMisAllZero}\,pp. The four misspecification configurations also belong to the primary matrix, so the correctly
specified method was run twice on the same missions. The two runs differ only in the measured planning times, because
the stopping rule reads a wall clock. They gave identical utility in \val{JRepeatSameShare}\% of the missions, and the
mean utility of a method differed by up to \val{JRepeatDiffMax}\,pp between the runs. Differences of this size between
stopping rules should therefore not be read as systematic.

\subsection{Tier C: PX4 software-in-the-loop with Gazebo}\label{sec:res-tierc}
All \val{JNCRuns} held-out flights of tier C were completed, that is, \val{JNC} missions with four methods each. Two runs
did not reach their altitude within 90\,s and were repeated as specified. A further \val{JCNInvalid} runs had a real-time
factor between \val{JCRTFInvalidMin} and \val{JCRTFInvalidMax}, below the validity range of 1.6 to 2.2 stated in the
protocol. The protocol did not specify how to treat such runs, so they were repeated once with the same seed, like a
technical failure, and the original runs were kept. The repeated runs reached the same utility as the original ones, and
the comparisons below are identical with either set (\val{JCDFixedFortyOrig}\,pp and \val{JCDASOrig}\,pp with the
original runs). The real-time factor of the analysed runs was \val{JCRTF} on average and at least \val{JCRTFmin}. A
released plan reached all vehicles after \val{JCLat}\,s of simulation time on average, which is included in the expected
latency of AnytimeSwarm-R.

Table~\ref{tab:tierc} summarises the results. AnytimeSwarm-R reached a utility of \val{JCURisk}\%, against
\val{JCUFixedForty}\% for ALNS with 40\,ms, \val{JCUAS}\% for AnytimeSwarm and \val{JCUFixed}\% for the full budget.
The paired differences were \val{JCDFixedForty}\,pp to ALNS with 40\,ms and \val{JCDAS}\,pp to AnytimeSwarm, with
Holm-adjusted $p$-values of \val{JCPFixedForty} and \val{JCPAS}. The fourth hypothesis is therefore not supported. With
\val{JNC} missions, the intervals are much wider than in tier B and cannot resolve differences of the size seen there.
Two effects of tier B reappear in tier C. The full budget again had the lowest mean utility, although its interval
includes zero (\val{JCDFixed}\,pp to AnytimeSwarm-R), and it was the only method that lost a noticeable number of tasks
to late plans,
\val{JCDropFixed} per mission against at most \val{JCDropAS} for the others. AnytimeSwarm-R again computed least among
the adaptive rules, with \val{JCMsRisk}\,ms per call against \val{JCMsAS}\,ms for AnytimeSwarm, a reduction of
\val{JCMsRedAS}\% (\val{JCDMsAS}\,ms).

\inputtable{tab_tierc}

The same \val{JNC} missions were also run in tier B with the rescaled parameters of tier C. Tier B predicted the outcome
of a flight well. The utility of a run in tier C correlated with that of the same run in tier B with $r=$ \val{JBCr}
(\val{JBCn} runs), and the mean utility per mission with $r=$ \val{JBCrMission} (Fig.~\ref{fig:bc}b). Tier C served
slightly more reward on average, \val{JBCUC}\% against \val{JBCUB}\%, so the kinematic model of tier B was, if anything,
slightly pessimistic for these missions. The small
differences between the methods were not predicted as well. The method means of the two tiers had a Kendall rank
correlation of \val{JBCTau}, and the sign of the per-mission difference between AnytimeSwarm-R and ALNS with 40\,ms
agreed in \val{JBCSignAgreeFixedForty}\% of the missions, counting ties as agreement only when both differences are zero.
Part of the per-run correlation is by construction, because both tiers draw the wind, speed and service disturbances
from the same random numbers, and much of it reflects the different difficulty of the missions. The correlation was
first computed with an erroneous pairing of runs in the analysis script, which was corrected after the results had
been seen; the hypothesis tests were not affected. In tier B, the advantage of AnytimeSwarm-R on these missions
was \val{JBgzDFixedForty}\,pp over ALNS with 40\,ms and \val{JBgzDAS}\,pp over AnytimeSwarm (Fig.~\ref{fig:bc}a). Both
tiers thus agree that the methods differ by at most a few percentage points on this scale, and that the full budget is
the weakest search method, while the order of the three short-computation methods is within the noise of both tiers.

\begin{figure}[t]
\centering
\fig{fig_bc}{\linewidth}
\caption{Tier B against tier C on the same \val{JNC} held-out missions (three UAVs, mid disturbances, $\kappa_B=50$):
(a) paired difference of AnytimeSwarm-R minus each baseline with 95\% bootstrap intervals; (b) utility of every run in
tier C against the same run in tier B (dashed: equality).}\label{fig:bc}
\end{figure}


\section{Study 2: regimes and software-in-the-loop details}\label{sec:Sstudy2}
Fig.~\ref{fig:Ss2main} shows the paired differences of Table~\ref{tab:s2main} of the main paper together with the
computation of every method.

\begin{figure}[t]
\centering
\fig{fig_s2_main}{\linewidth}
\caption{Study~2, tier B, held-out seeds 5000--5029: (a) utility of the proposed replanner minus each method, pooled
over \val{KNConfigs} configurations, with 95\% paired bootstrap intervals over \val{KNMissions} missions; (b) mean
utility against mean computation per call. Ablations are omitted.}\label{fig:Ss2main}
\end{figure}

\subsection{Regimes}\label{sec:s2-regime}
Table~\ref{tab:s2regime} and Fig.~\ref{fig:s2regime} break the utility down by disturbance level, time-window factor and
latency factor. Per configuration and without correction for multiple testing, the proposed replanner was better than
ALNS with 40\,ms in \val{KWinForty} of the \val{KNConfigs} configurations and worse in \val{KLossForty}, better than the
constant buffer in \val{KWinBuf} and worse in \val{KLossBuf}, and better than the conventional planner in \val{KWinConv}
and worse in \val{KLossConv}. The advantage over ALNS with 40\,ms was \val{KDFortyKTen}\,pp, \val{KDFortyKFifty}\,pp and
\val{KDFortyKTwoHundred}\,pp for $\kappa_B=10$, 50 and 200, and the advantage over online prediction stayed between
\val{KDOnlKTen}\,pp and \val{KDOnlKTwoHundred}\,pp.

The advantage was largest with the original windows. Under high disturbances, the proposed replanner served
\val{KDFortyWOneHigh}\,pp more reward than ALNS with 40\,ms and \val{KDBufWOneHigh}\,pp more than the constant buffer.
The losses all occurred under high disturbances with halved windows. In this regime, the proposed replanner reached
\val{KUPropWHalfHigh}\%, \val{KDBufWHalfHigh}\,pp relative to the constant buffer and \val{KDFortyWHalfHigh}\,pp relative
to ALNS with 40\,ms. It served \val{KServedPropWHalfHigh} tasks per mission, of which \val{KLatePropWHalfHigh} late,
against \val{KServedBufWHalfHigh} and \val{KLateBufWHalfHigh} for the constant buffer. The chance-constrained model was
thus too cautious when the windows were short compared with the predicted finish-time error. The effective deadline of
many tasks then probably falls before the earliest possible arrival, so the planner gives up tasks that a less cautious
plan still serves in time. The value of $z$ was selected per disturbance level on development missions, averaged over
both window widths. Selecting it also per window width, or adapting it online, is a natural extension.

\inputtable{tab_s2_regime}

\begin{figure}[t]
\centering
\fig{fig_s2_regime}{\linewidth}
\caption{Study~2, tier B: utility per regime as a function of the latency factor $\kappa_B$, for mid and high
disturbances with the original and with halved time windows (mean over \val{KNMissions} held-out missions per point).
The proposed replanner is shown with stars.}\label{fig:s2regime}
\end{figure}

\subsection{Tier C: PX4 software-in-the-loop with Gazebo}\label{sec:Ss2-tierc}
All \val{KCRuns} held-out flights of study~2 were completed, that is, \val{KCN} missions with four methods each. One run
failed to start and was repeated. As specified in the protocol, \val{KCInvalid} runs with a real-time factor outside
1.6 to 2.2 were repeated once with the same seed, and the original runs were kept. The laptop was also used for other
work during the campaign, and many repeats ran slower than the original runs. Of the \val{KCRuns} analysed runs,
\val{KCValid} lie in the validity range. The real-time factor was \val{KCRTF} on average, with quartiles of
\val{KCRTFq} and a minimum of \val{KCRTFmin}. Two sensitivity analyses that were not pre-registered therefore repeat the
comparisons with the original runs instead of the repeats, and with only those pairs of runs that both lie in the
validity range.

Table~\ref{tab:s2tierc} summarises the results. The proposed replanner reached a utility of \val{KCUProp}\%, against
\val{KCUForty}\% for ALNS with 40\,ms, \val{KCUBuf}\% for the constant buffer and \val{KCUConv}\% for the conventional
planner. The paired differences were \val{KCDForty}\,pp to ALNS with 40\,ms and \val{KCDConv}\,pp to the conventional
planner, with Holm-adjusted $p$-values of \val{KCPForty} and \val{KCPConv}. The fourth hypothesis is therefore not
supported. The secondary comparison with the constant buffer gave \val{KCDBuf}\,pp. The sensitivity analyses lead to
the same conclusion. With the original runs, the differences were \val{KCDFortyOrig}\,pp and \val{KCDConvOrig}\,pp.
With valid runs only, they were \val{KCDFortyValid}\,pp on \val{KCNFortyValid} missions and \val{KCDConvValid}\,pp on
\val{KCNConvValid} missions. No interval in the three analyses excludes zero.

\inputtable{tab_s2_tierc}

The mechanism seen in tier B reappears in the counts of late and skipped tasks. The proposed replanner served
\val{KCLateProp} tasks late per mission, against \val{KCLateForty} for ALNS with 40\,ms, which is \val{KCLateRedForty}\%
fewer. It skipped no task at departure, whereas ALNS with 40\,ms skipped \val{KCSkipForty} and the conventional planner
\val{KCSkipConv} tasks per mission. It computed \val{KCMsProp}\,ms per call, against \val{KCMsForty}\,ms for the fixed
budget. The conventional planner called the search \val{KCCallsConv} times per mission instead of about
\val{KCCallsProp}, as in tier B.


The same \val{KCN} missions were also run in tier B with the rescaled parameters of tier C. As in study~1, tier B
predicted the utility of a flight closely, with a correlation of $r=$ \val{KBCr} over \val{KBCn} runs. On these
missions, however, even tier B did not resolve the differences between the methods (Fig.~\ref{fig:s2bc}). Its estimates
were \val{KBgzDForty}\,pp to ALNS with 40\,ms, \val{KBgzDBuf}\,pp to the constant buffer and \val{KBgzDConv}\,pp to the
conventional planner. The tier C intervals overlap them. The per-mission difference to ALNS with 40\,ms had a standard
deviation of \val{KCSdForty}\,pp in tier C. Confirming the tier B estimate of \val{KCEffForty}\,pp on these missions with
a power of 80\% would require about \val{KCNPowerForty} missions per method. Tier C thus shows that the proposed
replanner runs on an autopilot and a physics engine with the same computation and fewer late tasks, but it cannot
confirm the size of the gain measured in tier B.

\begin{figure}[t]
\centering
\fig{fig_s2_bc}{0.62\linewidth}
\caption{Study~2, tier B against tier C on the same \val{KCN} held-out missions (three UAVs, mid disturbances,
$\kappa_B=50$): paired difference of the proposed replanner minus each baseline with 95\% bootstrap intervals.}
\label{fig:s2bc}
\end{figure}


\bibliographystyle{elsarticle-num}
\bibliography{references}
\end{document}
